\documentclass[10pt,a4paper]{article}

% ----------------- Paquetes -------------------%
\usepackage[T1]{fontenc}
\usepackage[utf8]{inputenc}
\usepackage[a4paper,margin=2.5cm]{geometry}
\usepackage{mathptmx}             
\usepackage{changepage} 
\usepackage{setspace}
\usepackage[spanish]{babel}
\usepackage[labelfont=bf,labelsep=space]{caption}
\usepackage{amssymb}
\usepackage{amsmath}
\usepackage{ebgaramond}
\usepackage{placeins}
\usepackage{etoolbox}
\usepackage{setspace}
\usepackage{parskip} 
\usepackage{tcolorbox}
\usepackage{needspace}
\usepackage{xparse}
\usepackage{ocgx2}
\usepackage{hyperref}
\usepackage{hypcap}
\usepackage{endnotes}
\usepackage{soul}\sethlcolor{yellow}% resaltado amarillo: \hl{texto} (Panel TCEE, ⌃H)


%%%%%%%%%%%%%%%%%%%%%%%%%%%%%%%%%%%%%%%%%%%%%%%%%%%%%%%%%%%%%%%%
% --- Fija primero tus valores globales "buenos" ---
\setstretch{1.2}                 % Interlineado global
\setlength{\parskip}{0.7\baselineskip}
\setlength{\parindent}{0pt}
\newlength{\GlobalParskip}
\setlength{\GlobalParskip}{\parskip}

\newcounter{eqnum}[subsection]
\renewcommand{\theeqnum}{\arabic{eqnum}}
\newcounter{alphaitem}
\newcounter{numitem}

%% ------------------- Recuadros----------------------%%
\newcommand{\puntosclave}[1]{%
  \begin{tcolorbox}[
    colframe=black,
    title={Puntos clave},
    before upper={\setlength{\parindent}{0.5cm}\setlength{\parskip}{0.5\baselineskip}}
  ]
  #1
  \end{tcolorbox}
}

\newcommand{\modificaciones}[1]{
  \begin{tcolorbox}[
    colback=white,
    colframe=red,
    title={Modificaciones a realizar},
    before upper={\setlength{\parindent}{0.5cm}\setlength{\parskip}{0.5\baselineskip}}
  ]
  #1
  \end{tcolorbox}
}

%% ------------------- Preguntas TEST----------------------%%
% Opciones: radio (single-choice)
\NewDocumentCommand{\OptRadio}{m m m}{%
  \noindent\CheckBox[radio,name=#1,value=#2,width=1.2ex,height=1.2ex]{}%
  \;\textbf{#2.}~#3\par
}

% Opciones: checkbox (multi-choice)
\NewDocumentCommand{\OptCheck}{m m}{%
  \noindent\CheckBox[name=#1,width=1.2ex,height=1.2ex,bordercolor={0 0 0}]{}%
  \;#2\par
}

% Pregunta single-choice (radio)
% Args: 1=id, 2=enunciado, 3=opciones, 4=solución+justificación, 5=imagen (o {}), 6=origen
\NewDocumentCommand{\PreguntaTestSingle}{m m m m m m}{%
  \par\medskip
  \noindent\textbf{#1.}~#2\par\smallskip

    % Imagen (si no es {})
    \IfBlankTF{#5}{}{%
      \begin{center}
        \includegraphics[width=0.75\linewidth]{#5}
      \end{center}
    }

    \begin{Form}
      #3
    \end{Form}

    \par\smallskip
    \switchocg{sol-#1}{\fbox{Mostrar / ocultar solución}}%
    \hfill{\footnotesize #6}\par\smallskip

    \begin{ocg}{Solución}{sol-#1}{0}
      \noindent #4
    \end{ocg}
  \par\medskip
}

% Pregunta multi-choice (checkboxes)
\NewDocumentCommand{\PreguntaTestMulti}{m m m m m m}{%
  \par\medskip
  \noindent\textbf{#1.}~#2\par\smallskip

    % Imagen (si no es {})
    \IfBlankTF{#5}{}{%
      \begin{center}
        \includegraphics[width=0.75\linewidth]{#5}
      \end{center}
    }
    \begin{Form}
      #3
    \end{Form}

    \par\smallskip
    \switchocg{sol-#1}{\fbox{Mostrar / ocultar solución}}%
    \hfill{\footnotesize #6}\par\smallskip

    \begin{ocg}{Solución}{sol-#1}{0}
      \noindent #4
    \end{ocg}
  \par\medskip
}

%% ------------------- Listas----------------------%%
    % --- Lista alfabética ---
    \newenvironment{la}
      {\begin{list}{\alph{alphaitem})}{%
          \usecounter{alphaitem}%
          \setlength{\leftmargin}{1cm}%
          \setlength{\parskip}{3pt}% 
          \setlength{\itemsep}{0pt}% ← espacio entre ítems
          \setlength{\parsep}{0pt}%  ← párrafos dentro de un ítem
      }}
      {\end{list}}
      
    % --- Lista numérica ---
    \newenvironment{lnum}
      {\begin{list}{\arabic{numitem}.}{%
          \usecounter{numitem}%
          \setlength{\leftmargin}{1cm}%
          \setlength{\parskip}{3pt}% 
          \setlength{\itemsep}{0pt}% ← espacio entre ítems
          \setlength{\parsep}{0pt}%  ← párrafos dentro de un ítem
      }}
      {\end{list}}

%% ------------------- Preguntas Test ----------------------%%
      \usepackage{xparse} % si ya lo tienes, no lo repitas

    \NewDocumentEnvironment{test}{m m o}{%
      \begin{tcolorbox}[
        colback=white,
        colframe=black!15,
        boxrule=0.6pt,
        arc=2mm,
        left=2mm,right=2mm,top=2mm,bottom=2mm
      ]
      \centering
      \includegraphics[width=\linewidth]{#1}
      \end{tcolorbox}
    
      \vspace{2mm}
    
      % Caja SOLO para texto (SÍ partible y mantiene márgenes al partir)
      \begin{tcolorbox}[
        enhanced,
        breakable,
        colback=black!3,
        colframe=black!25,
        boxrule=0.6pt,
        arc=2mm,
        left=2mm,right=2mm,top=1.5mm,bottom=1.5mm
      ]
      \textbf{Respuesta:} \textbf{#2}%
      \IfValueT{#3}{\par\smallskip \textbf{Justificación:} #3}
      \end{tcolorbox}
    }{}
      
% ------------------- Imágenes----------------------%
    \graphicspath{{Img/}}% Rutas por defecto 
    % --- Imagen adaptativa ---
    % Registros de dimensión definidos una sola vez
    \newdimen\imgcap
    \newdimen\imgmin
    \newdimen\imgmax
    \newdimen\imgremain
    \newdimen\imgh
    \newdimen\imgneed
    
    \newcommand{\imagenfit}[3]{%
      \addvspace{10pt}% separación antes
      \par\begingroup
        % Parámetros de composición (asignaciones, no \newdimen)
        \imgcap=1\baselineskip            % alto estimado de título+fuente
        \imgmin=0.15\textheight
        \imgmax=0.15\textheight
        \imgremain=\pagegoal
        \ifdim\pagegoal=\maxdimen \imgremain=0pt\fi
        \advance\imgremain by -\pagetotal
        \advance\imgremain by -\imgcap
        % Decisión de altura destino
        \ifdim\imgremain<\imgmin
          \newpage
          \imgh=0.20\textheight
        \else
          \imgh=\imgremain
          \ifdim\imgh>\imgmax \imgh=\imgmax \fi
        \fi
        % Reservar espacio para evitar cortes del bloque completo
        \imgneed=\imgh \advance\imgneed by \imgcap
        \Needspace{\imgneed}
        % Cuerpo
        \noindent\begin{minipage}{\textwidth}
          \centering
          \captionof{figure}{\textemdash\ #2}\par\vspace{0.3em}% título arriba
          \includegraphics[width=\linewidth,height=\imgh,keepaspectratio]{\detokenize{#1}}\par
          \vspace{0.3em}\small\textit{Fuente: }#3% fuente abajo
        \end{minipage}%
      \endgroup
      \par\addvspace{10pt}%
    }

%------------ Bloque de RECUADRO ESPECIAL -------------%
    \newenvironment{specialblock}[1]{%
      \begin{quote} % crea sangría a izquierda y derecha
      \small % texto más pequeño
      \noindent\textbf{#1}\\[-0.3em] % título en negrita
      \rule{\linewidth}{0.4pt}\\[0.6em]}{
      \end{quote}}
      
%-------------------------- Portada -----------------------%
\newcommand{\oldtitle}[1]{\def\OldTitle{#1}}
\newcommand{\changerationale}[1]{\def\ChangeWhy{#1}}

\newcommand{\changebox}{%
\begin{tcolorbox}[title={Historial de cambios del tema}]
\textbf{Título anterior:} \OldTitle\par
\textbf{Motivo del cambio:} \ChangeWhy
\end{tcolorbox}
}

%-------------------------- Author Font -----------------------%
    \newcommand{\authorfont}[1]{{\fontfamily{EBGaramond-TLF}\selectfont #1}}

%-------------------------- Anexos -----------------------%
    \makeatletter
    \newcounter{anexo}
    \renewcommand{\theanexo}{\Roman{anexo}}
    \newcommand{\anexo}[1]{%
      \clearpage
      \phantomsection
      \refstepcounter{anexo}%
      \noindent\textbf{Anexo \theanexo.- }#1\par\medskip
      \addcontentsline{stc}{section}{Anexo \theanexo.- #1}%
    }
    \makeatother

    %-------------------------- Lista de figuras (personal) -----------------------%
\makeatletter
\newcommand{\MyListOfFigures}{%
  \clearpage
  \phantomsection
  {\centering\bfseries\Large Índice de figuras\par}%
  \medskip
  % Entrada en nuestro índice de contenido (stc)
  \addcontentsline{stc}{section}{Índice de figuras}%
  % Recuperación del listado (sin cabecera propia)
  \begingroup
    \parindent=0pt\parskip=2pt
    \@starttoc{lof}%
  \endgroup
}
\makeatother

%-------------------------- Bibliografía -----------------------%
\makeatletter
% Contador para las referencias
\newcounter{myref}
% Cabecera de la bibliografía, entra en el índice 'stc'
\newcommand{\MyBibliography}{%
  \clearpage
  \phantomsection
  {\centering\bfseries\Large Bibliografía y referencias\par}%
  \medskip
  \addcontentsline{stc}{section}{Bibliografía y referencias}%
  \setcounter{myref}{0}%
}
\newcommand{\MyBibItem}[4]{%
  \refstepcounter{myref}%
  \noindent[\themyref]\ #1.\ \emph{#2}. #3. #4\par
}
\makeatother

%--------- Establecer cuatro niveles de índice --------%
    \setcounter{tocdepth}{4}   
    \setcounter{secnumdepth}{4}% numera hasta \paragraph

%%%%%%%%%%%%%%%%%%%%%%%%%%%%%%%%

\makeatletter

    %--------- Interlineado Global --------%
    \edef\GlobalBaseLineStretch{\baselinestretch}
    
    %--------- Espaciado de seccinones --------%
    \renewcommand\section{\@startsection{section}
      {1}{\z@}
      {2.5ex \@plus 1ex \@minus .2ex} % espacio antes
      {1.5ex \@plus .2ex}             % espacio después
      {\normalfont\Large\bfseries}    % formato del título
    }
    \renewcommand\subsection{\@startsection{subsection}{2}{\z@}
      {3ex \@plus 0.8ex \@minus .2ex}
      {1.5ex \@plus .2ex}
      {\normalfont\large\bfseries}
    }
    \renewcommand\subsubsection{\@startsection{subsubsection}{3}{\z@}
      {3ex \@plus 0.5ex \@minus .2ex}
      {1ex \@plus .2ex}
      {\normalfont\normalsize\bfseries}
    }
    \renewcommand\paragraph{\@startsection{paragraph}{4}{\z@}%
    {-2ex \@plus -0.5ex \@minus -.2ex}% espacio antes
    {0.8ex \@plus .2ex}% espacio después
    {\normalfont\normalsize\itshape}}% ← cursiva en lugar de negrita

    %--------- Ecuaciones --------%   
    % Variables globales para almacenar títulos y notas
    \gdef\leq@current@section{}%
    \gdef\leq@current@subsection{}%
    \gdef\leq@last@printed@section{}%
    \gdef\leq@last@printed@subsection{}%
    
    % Comando para añadir nota (invisible en el cuerpo)
    \newcommand{\eqnote}[1]{%
      \addcontentsline{leq}{leqnote}{#1}%
    }
    
    % Comando para ecuaciones
    \newcommand{\eqblock}[2]{%
      % Verificar si necesitamos imprimir el título de sección
      \ifx\leq@current@section\leq@last@printed@section\else
        \addcontentsline{leq}{leqsection}{\leq@current@section}%
        \global\let\leq@last@printed@section\leq@current@section
        \gdef\leq@last@printed@subsection{}% Reset subsección al cambiar sección
      \fi
      % Verificar si necesitamos imprimir el título de subsección
      \ifx\leq@current@subsection\leq@last@printed@subsection\else
        \ifx\leq@current@subsection\@empty\else
          \addcontentsline{leq}{leqsubsection}{\leq@current@subsection}%
          \global\let\leq@last@printed@subsection\leq@current@subsection
        \fi
      \fi
      % Ecuación
      \refstepcounter{eqnum}%
      \begin{equation*}
        #1\tag{E.\theeqnum}
      \end{equation*}%
      \addcontentsline{leq}{leqt}{[E.\theeqnum] -- #2}%
      \addcontentsline{leq}{leqm}{#1}%
    }
    
    %%% ----- Índice de ecuaciones ----------%%%
    % Tamaño para []
    \newcommand{\leqIndexFont}{\normalsize}
    \newcommand{\leqTitleFont}{\tiny}
    \newcommand{\leqNoteFont}{\tiny}
    % Cabecera de sección 
    \newcommand*\l@leqsection[2]{%
      \addvspace{25pt}% Mayor separación antes de nueva sección
      \noindent\leqIndexFont
      \makebox[\linewidth]{\rule{.38\linewidth}{0.4pt}\ \ \textbf{  \MakeUppercase{#1}}\ \ \rule{.38\linewidth}{0.4pt}}%
      \par\vspace{-10pt}
    }
    % Cabecera de subsección 
    \newcommand*\l@leqsubsection[2]{%
      \par\addvspace{15pt}%
      \noindent\leqIndexFont #1
      \par\vspace{-7pt}
    }
    % Título de ecuación 
    \newcommand*\l@leqt[2]{%
    \par\addvspace{10pt}%
      \par\noindent\leqTitleFont #1\par
    }
    % Miniatura de ecuación
    \newcommand*\l@leqm[2]{%
      \vspace{0.3\baselineskip}%
      \par\scriptsize\noindent\hspace*{1.5em}\(#1\)\par%
    }
    % Nota de ecuación 
    \newcommand*\l@leqnote[2]{%
      \vspace{0.2\baselineskip}%
      \par\noindent\leqNoteFont\hspace*{1.5em}\textbf{Nota.--} #1\par\vspace{0.5\baselineskip}%
    }
    % Generación del índice
    \newcommand{\mylistofequations}{%
      \clearpage
      \phantomsection
      % Título visual de la lista (ajústalo a tu gusto):
      {\centering\bfseries\Large Índice de ecuaciones\par}
      % Entrada en nuestro índice 'stc':
      \addcontentsline{stc}{section}{Índice de ecuaciones}%
      % Llamada a tu lista real (usa el nombre exacto de tu macro):
      \@starttoc{leq}%
    }
    
    %--------- Captura de títulos de sección --------%
    \let\leq@original@section\section
    \let\leq@original@subsection\subsection
    % Flag para saber si estamos en una sección asterisco
    \newif\ifLeqStarSection
    \LeqStarSectionfalse
    
    % Redefinir \section CON CONTROL DE ASTERISCO
    \renewcommand{\section}{%
      \@ifstar{\leq@section@star}{\@dblarg\leq@section@nostar}%
    }
    \def\leq@section@star#1{%
      \global\LeqStarSectiontrue
      \gdef\leq@current@section{#1}%
      \gdef\leq@current@subsection{}%
      \leq@original@section*{#1}%
      \global\LeqStarSectionfalse
    }
    \def\leq@section@nostar[#1]#2{%
      \global\LeqStarSectionfalse
      \gdef\leq@current@section{#2}%
      \gdef\leq@current@subsection{}%
      \leq@original@section[#1]{#2}%
    }
    % Redefinir \subsection
    \renewcommand{\subsection}{%
      \@ifstar{\leq@subsection@star}{\@dblarg\leq@subsection@nostar}%
    }
    \def\leq@subsection@star#1{%
      \global\LeqStarSectiontrue
      \gdef\leq@current@subsection{#1}%
      \leq@original@subsection*{#1}%
      \global\LeqStarSectionfalse
    }
    \def\leq@subsection@nostar[#1]#2{%
      \global\LeqStarSectionfalse
      \gdef\leq@current@subsection{#2}%
      \leq@original@subsection[#1]{#2}%
    }

    % ----- Restauración de valores Globales de estilo ----------%
    % Flags
    \newif\ifInFootnote
    \InFootnotefalse
    \pretocmd{\@footnotetext}{\InFootnotetrue}{}{}
    \apptocmd{\@footnotetext}{\InFootnotefalse}{}{}
    % Profundidad de listas (soporta anidadas)
    \newcount\MyListDepth \MyListDepth=0
    \AtBeginEnvironment{la}{\advance\MyListDepth by 1}
    \AfterEndEnvironment{la}{\advance\MyListDepth by -1}
    \AtBeginEnvironment{lnum}{\advance\MyListDepth by 1}
    \AfterEndEnvironment{lnum}{\advance\MyListDepth by -1}
    \AtBeginEnvironment{itemize}{\advance\MyListDepth by 1}
    \AfterEndEnvironment{itemize}{\advance\MyListDepth by -1}
    \AtBeginEnvironment{enumerate}{\advance\MyListDepth by 1}
    \AfterEndEnvironment{enumerate}{\advance\MyListDepth by -1}
    % Restauración diferida: hazlo al INICIO del próximo párrafo real
    \newif\if@pendingRestore
    \@pendingRestorefalse
    \newcommand{\RequestRestore}{\global\@pendingRestoretrue}
    \everypar{%
      \if@pendingRestore
        \global\@pendingRestorefalse
        \ifInFootnote\else
          \ifnum\MyListDepth=0
            \setlength{\parskip}{\GlobalParskip}%
            \setstretch{\GlobalBaseLineStretch}%
          \fi
        \fi
      \fi
    }
    % Al final de bloques:
    \AfterEndEnvironment{equation}{\RequestRestore}
    \AfterEndEnvironment{equation*}{\RequestRestore}
    \AfterEndEnvironment{la}{\RequestRestore}
    \AfterEndEnvironment{lnum}{\RequestRestore}
    % Al final de macros:
    \apptocmd{\eqblock}{\RequestRestore}{}{}
    \apptocmd{\imagenfit}{\RequestRestore}{}{}
    
\makeatother

% ===================== ToC propio con 4 niveles (único bloque) =====================
\makeatletter

% --- Escritores: añadimos una línea al .stc en cada cabecera numerada ---
% Guardamos las versiones actuales para llamar al comportamiento normal
\let\MyTOC@orig@section\section
\let\MyTOC@orig@subsection\subsection
\let\MyTOC@orig@subsubsection\subsubsection
\let\MyTOC@orig@paragraph\paragraph

% SECTION
\renewcommand{\section}{%
  \@ifstar{\MyTOC@section@star}{\@dblarg\MyTOC@section@nostar}%
}
\def\MyTOC@section@star#1{%
  \MyTOC@orig@section*{#1}% sin entrada al .stc
}
\def\MyTOC@section@nostar[#1]#2{%
  \MyTOC@orig@section[{#1}]{#2}% comportamiento normal (escribe al .toc)
  % Escribimos también nuestra línea al .stc (texto con \numberline)
  \addcontentsline{stc}{section}{\protect\numberline{\thesection}#2}%
}

% SUBSECTION
\renewcommand{\subsection}{%
  \@ifstar{\MyTOC@subsection@star}{\@dblarg\MyTOC@subsection@nostar}%
}
\def\MyTOC@subsection@star#1{%
  \MyTOC@orig@subsection*{#1}%
}
\def\MyTOC@subsection@nostar[#1]#2{%
  \MyTOC@orig@subsection[{#1}]{#2}%
  \addcontentsline{stc}{subsection}{\protect\numberline{\thesubsection}#2}%
}

% SUBSUBSECTION
\renewcommand{\subsubsection}{%
  \@ifstar{\MyTOC@subsubsection@star}{\@dblarg\MyTOC@subsubsection@nostar}%
}
\def\MyTOC@subsubsection@star#1{%
  \MyTOC@orig@subsubsection*{#1}%
}
\def\MyTOC@subsubsection@nostar[#1]#2{%
  \MyTOC@orig@subsubsection[{#1}]{#2}%
  \addcontentsline{stc}{subsubsection}{\protect\numberline{\thesubsubsection}#2}%
}

% PARAGRAPH
\renewcommand{\paragraph}{%
  \@ifstar{\MyTOC@paragraph@star}{\@dblarg\MyTOC@paragraph@nostar}%
}
\def\MyTOC@paragraph@star#1{%
  \MyTOC@orig@paragraph*{#1}%
}
\def\MyTOC@paragraph@nostar[#1]#2{%
  \MyTOC@orig@paragraph[{#1}]{#2}%
  % Si no numeras los \paragraph, cambia \theparagraph por texto plano sin \numberline
  \addcontentsline{stc}{paragraph}{\protect\numberline{\theparagraph}#2}%
}

% --- Impresor del índice propio (lee el .stc) ---
\newcommand{\MyTOC}{%
  \par\bigskip
  {\centering\bfseries\Large Índice de contenido\par}%
  \medskip
  \begingroup
    \parindent=0pt\parskip=2pt
    \@starttoc{stc}%
  \endgroup
}

\makeatother
% ===================== fin ToC propio =====================


%%%%%%%%%%%%%%


% --- Datos del tema ---
\title{Tema 3.A.10 -- Teoría de la demanda del consumidor (III)\\
\large Elección del consumidor en situaciones de riesgo e incertidumbre}
\author{Marcelo Ortega}
\date{Febrero 2026}
\newcommand{\TemaCode}{3A10}

\oldtitle{Tema 3.A.7. Teoría de la elección del consumidor en situaciones de riesgo e incertidumbre}
\changerationale{
    Sin cambios
}

\begin{document}


\maketitle
\changebox

\newpage

% --- Página de resumen y modificaciones ---
\puntosclave{ %% Escribir puntos clave Apuntes CECO
    }
\vspace{1cm}

\modificaciones{
    
    }

\newpage
\MyTOC
\newpage

\mylistofequations
\clearpage

\MyListOfFigures
\clearpage


\pagenumbering{arabic}
\setcounter{page}{1}


\clearpage
\section{Introducción}\label{intro}
ALFRED MARSHALL, en sus Principios de Economía (1890) define la economía como \textit{la ciencia de la vida diaria en lo que respecta a las acciones humanas tomadas para alcanzar un nivel máximo de bienestar.}

Esta definición nos muestra cómo uno de los principios subyacentes a la reflexión económica, pero particularmente enfatizado en la teoría neoclásica, es el del individualismo metodológico\footnote{%
    El individualismo metodológico es un método ampliamente utilizado en las ciencias sociales. Sostiene que todos los fenómenos sociales —estructura y cambios— son en principio explicables por elementos individuales, es decir, por las propiedades de los individuos, como pueden ser sus metas, sus creencias y sus acciones. Sus defensores lo ven como una filosofía-método destinada a la explicación y comprensión amplia de la evolución de toda la sociedad como el agregado de las decisiones de los particulares. En principio es un reduccionismo, es decir, una reducción de la explicación de todas las grandes entidades con referencias en las más pequeñas.
}. Se contempla el objeto de la teoría como una realidad social compuesta de individuos que se interrelacionan en economías descentralizadas.

En su objetivo fundamental de comprender y predecir el funcionamiento de los mercados, la microeconomía examina el \textbf{comportamiento} de dos \textbf{agentes} fundamentales: consumidores y productores\footnote{%
    No hay que olvidar que la microeconomía contemporánea contempla esta separación estricta entre consumidores y productores como “una hipersimplificación del proceso por el que los bienes se compran y se consumen” (EKELUND y HÉBERT, 2013). Ejemplos que muestran el desdibujado de esta frontera son las “tecnologías del consumo”, es decir, la aplicación de la teoría de la producción a las decisiones de consumo, como son el enfoque de características de KEVIN LANCASTER, la economía doméstica de GARY BECKER, la producción doméstica de REUBEN GRONAU o la economía de la información de GEORGE J. STIGLER (la información sobre los bienes de consumo, como bien económico o costoso, obliga a un proceso de búsqueda que debe combinarse con el bien de consumo físico).
}.

\textbf{En la teoría de la demanda del consumidor, los individuos objeto de estudio son los consumidores}\footnote{%
    Los consumidores son aquellos agentes que actúan típicamente como demandantes de bienes en los mercados de productos y como oferentes en los mercados de factores productivos.
}. Se asume que estos se comportan de manera optimizadora y quedan \textbf{caracterizados por su deseo de consumir ciertos bienes sometidos a una restricción presupuestaria}.

El comportamiento se reduce en último término a la toma de decisiones entre alternativas. Estas alternativas conducen a resultados que pueden ser conocidos con certeza de antemano, o no.

Cuando no son conocidos hablamos de situaciones de riesgo o incertidumbre y la modelización de la decisión en estas condiciones tiene características peculiares. Esa incertidumbre respecto a los resultados de las decisiones puede ser objetiva o subjetiva:
\begin{la}
    \item Cuando la incertidumbre es objetiva se denomina \textit{riesgo}\footnote{%
        La palabra riesgo procede del veneciano “rischio”. Los “rischios” eran protuberancias rocosas a la entrada de los puertos que hundían barcos con una probabilidad estimable de antemano.
    }. Es objetiva, porque depende del objeto, no del sujeto: las probabilidades de cada estado de la naturaleza son perfectamente conocidas y dependen de la lotería en cuestión, no del agente decisor. 
    \item Por otro lado, siguiendo a KNIGHT (1921)\footnote{%
        KNIGHT (1921) propuso distinguir entre \textit{riesgo e incertidumbre} teniendo en cuenta si las probabilidades nos vienen dadas de manera objetiva o no.
    }, se encuentran los contextos de incertidumbre en los que las probabilidades asignadas a cada estado de naturaleza no vienen dadas de manera objetiva.
\end{la}

\subsection{Relevancia}\label{relevancia}
Podemos señalar 2 razones que justifican la importancia del estudio de la demanda del consumidor con información imperfecta:
\begin{lnum}
    \item Se trata de un desarrollo fundamental en la teoría de la demanda del consumidor, siendo la información imperfecta una característica ineludible de la vida económica.
    \item Los enfoques presentados en esta exposición pueden ser de tremenda utilidad en una enorme variedad de situaciones, no sólo a nivel teórico, pero también en muchos sectores relacionados con la incertidumbre (por ejemplo, decisiones de aseguramiento o de inversión).
\end{lnum}

\subsection{Contextualización}\label{contexto}
En los temas anteriores hemos tratado [Tema 8A y Tema 9A] los Problemas de la Elección (del consumidor) que suponían resultados perfectamente conocidos. Sin embargo, muchas decisiones importantes en economía llevan un elemento de \textit{riesgo} asociado. Aunque es formalmente posible analizar dichas situaciones utilizando la Teoría de la Elección general tratada en los temas anteriores, hay una buena razón para desarrollar una rama teórica más especializada: las alternativas inciertas tienen una estructura que podemos utilizar para restringir las preferencias que consumidores “racionales” puedan tener.

Desde una perspectiva histórica, es enormemente llamativo el origen del análisis económico de las decisiones en ausencia de perfecta certidumbre. El análisis de estas decisiones en ausencia de información perfecta se inicia de forma pionera por DANIEL BERNOULLI\footnote{%
    La hipótesis fue inicialmente planteada en 1738 a través del concepto de esperanza moral. Ésta afirmaba que las decisiones que toman los individuos no surgen solo de los resultados esperados sino también de otros factores definidos a través de una función de utilidad. Para Bernoulli, escoger una opción según su valor esperado no era racional porque en ese caso, dadas unas condiciones concretas (Paradoja de San Petesburgo, formulada por NICOLAUS BERNOULLI, dado un juego de apuesta cuyo valor esperado sea infinito) “lo racional” sería apostarlo todo. Curiosamente, todos estos miembros, como también: JACOB BERNOULLI, conocido por sus aportaciones al campo de la Probabilidad y Estadística (Distribución de Bernoulli), comparten relación de parentesco (todos son tíos, hermanos o sobrinos).
}, quien en el siglo XVIII realiza el primer análisis riguroso de la toma de decisiones en ausencia de certidumbre, \textbf{aplicando de forma adelantada el concepto de utilidad marginal} para dar respuesta a la conocida Paradoja de San Petersburgo\footnote{Formulada por NICOLAUS BERNOULLI por correspondencia con MONTMORT (1710-1712).} que no parecía resolverse aplicando el criterio de la ganancia esperada.

El análisis de las decisiones bajo incertidumbre quedaría relegado y sería rescatado ya entrado el siglo XX. Destacamos a FRANK KNIGHT, que en su obra “Risk, Uncertainty and Profit” (1921) fue pionero en distinguir entre riesgo e incertidumbre. Es habitual, sin embargo, utilizar el término incertidumbre para referirse a ambos fenómenos cuando no se pretende distinguir explícitamente.

En cualquier caso, no sería hasta los años 40, cuando JOHN VON NEUMANN y OSKAR MORGENSTERN publicaran su obra señera “The Theory of Games and Economic Behaviour” (1944)\footnote{8} que por primera vez establecía con rigor la derivación axiomática de la teoría de la utilidad esperada de DANIEL BERNOULLI.

En los años 50 tendría lugar el segundo gran hito. Sería la obra de LEONARD SAVAGE \textit{“Foundations of Statistics”} (1954), que consiguió derivar la utilidad esperada sin necesidad de imponer la existencia de probabilidades objetivas (i.e. en situaciones de incertidumbre).

Más recientemente, también han surgido voces discordantes que ponen de manifiesto las limitaciones de los análisis dominantes y con ello han promovido la mejora y el desarrollo de este campo. Entre otros caben destacar las contribuciones de MAURICE ALLAIS\footnote{%
    MAURICE ALLAIS fue galardonado con el Premio Nobel de Economía en 1988 «Por sus contribuciones a la teoría de los mercados y la eficiente utilización de los recursos».
}, DANIEL KAHNEMAN\footnote{%
    DANIEL KAHNEMAN fue galardonado con el Premio Nobel de Economía en 2002 junto con VERNON SMITH «Por integrar aspectos de la teoría psicológica sobre el comportamiento económico del ser humano en momentos de incertidumbre y realizar análisis empíricos de laboratorio, especialmente sobre mecanismos alternativos de mercado». La concesión del Premio Nobel de Economía del 2002 a DANIEL KAHNEMAN es un reconocimiento de la importancia de la psicología en la explicación del comportamiento económico. En colaboración con AMOS TVERSKY –quien, de no haber fallecido en 1996 hubiera recibido también el galardón– DANIEL KAHNEMAN formuló una nueva teoría de la decisión humana, la \textit{Prospect Theory}, que ha dado origen a una nueva rama del análisis financiero, la \textit{Behavioral Finance (Psicología de las finanzas).}
} o AMOS TVERSKY.

\subsection{Problemática (Preguntas clave)}\label{problematica}
\begin{lnum}
    \item \textit{¿Qué es el riesgo? ¿Cómo han sido modelizadas las decisiones en situaciones de riesgo?}
    \item \textit{¿Qué es la incertidumbre? ¿Cómo han sido modelizadas las decisiones en situaciones de incertidumbre?}
\end{lnum}
\section{Teoría de la Utilidad Esperada}
Siguiendo a KNIGHT podemos definir una situación de riesgo como aquella en la que la aleatoriedad a la que se enfrenta el agente económico se presenta en forma de probabilidades objetivas (especificadas exógenamente o calculables científicamente).

De este modo, la aleatoriedad depende del objeto, no del sujeto: \textbf{las probabilidades de cada estado de la naturaleza son perfectamente conocidas y dependen de la lotería en cuestión, no del agente decisor}.

Pueden ser consideradas una excesiva simplificación de la realidad, ya que es difícil encontrar ejemplos de riesgo puro en la actividad económica real. Es por ello que usamos ejemplos simples como tirar una moneda al aire, dados o ruletas de casinos.//
La Teoría de la Utilidad Esperada fue la primera aplicación del modelo de decisión neoclásico –que empezaba a concretarse en la primera mitad del siglo XX– a la elección bajo incertidumbre.

VON NEUMANN y MORGENSTERN desarrollaron los pilares en los que sostendrían posteriormente las finanzas modernas, el equilibrio general de ARROW-DEBREU o la modelización de fenómenos como la selección adversa o el riesgo moral.

\subsection{Ingredientes y obtención de la función de utilidad esperada}
Al igual que sucede con la elección del consumidor en condiciones de certidumbre [ver Tema 3.A.8], la metodología en este caso seguirá un esquema muy similar, en cuanto a que cuenta con las mismas etapas:
\begin{lnum}
    \item Se define un conjunto de elección posible;
    \item Sobre el que se establece la relación binaria de preferencias del consumidor; y, 
    \item Gracias a un conjunto de axiomas, podrá representarse mediante una función matemática concreta: la función de utilidad esperada. 
\end{lnum}
Esto nos permitirá aplicar las herramientas del cálculo diferencial y por tanto plantear y resolver un problema de maximización de la utilidad esperada sujeto a una restricción presupuestaria, del que se puede extraer una decisión óptima del consumidor en situaciones de riesgo.

En cualquier caso, el modelo de la utilidad esperada difiere de la teoría de la elección del consumidor en condiciones de certidumbre en dos aspectos:
\begin{la}
    \item Los objetos de elección ya no son cestas de consumo deterministas, sino loterías estocásticas, que, por lo tanto, dan lugar a aleatoriedad.
    \item El modelo de utilidad esperada impone una restricción muy específica en la forma funcional de la función de utilidad, al depender esta no solo de las preferencias del consumidor sino también de la distribución de probabilidad.
\end{la}
\subsubsection{Conjunto de elección: Loterías}
El ingrediente básico para la construcción de esta teoría es el concepto de lotería\footnote{%
    En una lotería simple, los resultados que pueden suceder son certeros, si sucede (por probabilidad) un elemento del conjunto, el resultado sociado es certero y son los integrantes (o integrante) de dicho elemento del conjunto general de loterías.
}, una herramiento formal que utilizaremos para representar alternativas con diversos grados de \textit{riesgo}.

[Lotería simple: Matemáticas]

En pocas palabras, una lotería es una distribución de probabilidad (discreta), representando la probabilidad de que se dé cada estado de la naturaleza o \textit{“resultado”} (equivalente).

Una variante más general de una \textit{lotería} permite que los resultados de una \textit{lotería} sean \textit{loterías} en sí mismos. Se trata verdaderamente de un proceso de elección en dos etapas.

[Lotería compuesta: Matemáticas]

El análisis teórico parte de un supuesto \textit{consecuencialista}: asumimos que para toda alternativa (incierta o riesgosa), solo la lotería sobre consecuencias finales es relevante para el sujeto. Esto es: no importa si la distribución de probabilidad sobre resultados finales se origina como resultado de una lotería simple o una lotería compuesta (por dos o más eventos)\footnote{%
    De hecho, nuestra hipótesis consecuencialista requiere que el agente vea estas dos loterías como equivalentes. Esto a veces se incluye como un axioma de las preferencias (axioma de equivalencia racional), según el individuo se mostrará indiferente entre una lotería compuesta y su lotería simple equivalente. Por ejemplo, el individuo se mostrará indiferente entre una lotería con una probabilidad del 25\% de ganar 100 € y otra lotería en la que tiene una probabilidad del 50\% de ganar un boleto de lotería que cuenta con un 50\% de ganar 100 €.
}. Para todo conjunto de loterías compuestas puede calcularse un correspondiente conjunto de loterías simples (o reducidas) que generan la misma distribución sobre resultados, esto es, el valor de probabilidad para cada elemento de la lotería simple (reducida) será el sumatorio de productos resultantes de multiplicar la probabilidad de que suceda una lotería determinada con la probabilidad de que, dentro de ésta, suceda el evento en cuestión\footnote{%
    Esta propiedad, aunque parezca banal, es de extrema importancia, pues permite demostrar la estructura lineal del espacio de loterías, propiedad central de la teoría de la elección en contextos de incertidumbre.
}. Por tanto, de ahora en adelante (por simplicidad) vamos a restringir el desarrollo al de una lotería simple (sin pérdida de generalidad). 
\subsubsection{Preferencia del consumidor sobre loterías y sus axiomas}
Habiendo desarrollado esta forma de modelar el riesgo entre alternativas, pasaremos a las preferencias que el agente tiene sobre éstas. 
\begin{itemize}
    \item Preferencias racionales (completitud y transitividad) permitiendo comparación entre cualquier par de elementos del conjunto.
    \item Independencia\footnote{}. [Nos permite imponer una estructura considerablemente mayor a esta Función de Utilidad (existente)]. En palabras, si mezclamos una combinación par de loterías con una tercera, entonces el orden de preferencias resultante de la “mezcla” no depende de la tercera lotería escogida (es independiente de…)
\end{itemize}
[Significado: Matemáticas]

El cumplimiento de estos tres axiomas garantiza que todas las loterías (completitud) pueden clasificarse en conjuntos de indiferencia (reflexividad), de modo que cada lotería pertenezca a una y solo una clase de indiferencia (transitividad).

Cabe destacar que, [MAS-COLELL, et. al, 1995], los axiomas de racionalidad impuestos en este contexto son más “restrictivas”, que lo que lo son en un contexto de información perfecta. Cuanto más complejas son las alternativas, mayor es la “carga” que sobre el modelo imponen los axiomas de racionalidad. De hecho, su “ajuste a la realidad” en un contexto de incertidumbre ha sido objeto de amplísimo debate.

Dos axiomas adicionales:
\begin{lnum}
    \item Continuidad. En palabras, este axioma implica que pequeños cambios en las probabilidades no cambian la naturaleza de la ordenación entre dos loterías. El axioma de continuidad implica la existencia de una función de utilidad que representa las preferencias del consumidor. 
\end{lnum}

\footnote{%
    Descripción 1: El Axioma de Independencia está en el centro de la teoría de la elección bajo condiciones de incertidumbre. Esto es precisamente por se sirve, de una manera fundamental, de la estructura de incertidumbre presente en este modelo. En la teoría neoclásica de demanda del consumidor [ver Tema 3.8.A y 3.9.B] no hay razón para pensar que las preferencias de un consumidor sobre determinadas cestas de bienes deban ser independientes de la cantidad de otros bienes que vaya a consumir. Sin embargo, en contextos de incertidumbre, \textbf{sí es natural} pensar que las preferencias de un consumidor entre dos loterías, pueda definir ex-ante cuál preferiría que formase parte de otra cesta de consumo (cualquiera), independientemente del resultado de esta otra lotería compuesta. Este “nuevo” resultado, no debería alterar sus preferencias puesto que, en contraste con el contexto del consumidor, no consume ninguna de las dos alternativas juntas con esta tercera opción, por el contrario, sólo en vez de esa. 
Descripción 2: La justificación de esta irrelevancia respecto de terceras alternativas descansa en que, a diferencia de la demanda de consumo cierto (donde esas terceras opciones (en ese caso cestas de bienes) pueden ser relevantes pues realmente se pueden combinar con otras en el consumo), en el ámbito del riesgo las loterías no pueden “consumirse” simultáneamente, sino que el agente se ve obligado a elegir necesariamente sólo una de ellas.
}
\footnote{%
    La teoría recibe su nombre de Utilidad Esperada porque en situaciones de riesgo se considera que el agente valora las diferentes loterías, no según la Utilidad de la Esperanza del resultado, \textbf{sino según la Esperanza de las Utilidades de los diferentes resultados}. Esto nos permite resolver junto a la idea de la utilidad marginal decreciente el juego de la paradoja de San Petersburgo (permitiendo explicar comportamientos observados). En el caso de este juego si actuáramos conforme a la utilidad de la esperanza, estaríamos dispuestos a pagar hasta una cantidad que nos proveyera de una utilidad de infinito (que era la esperanza matemática del resultado). Según este enfoque, será la esperanza de la utilidad de que se produzca cada uno de los resultados. Nota personal: Cuando pienses en la utilidad esperada a nivel analítico, pensar en inglés “Expected Utility” ya que se ve más claro que estamos hablando de la esperanza de la utilidad: $E\big[u(\tilde{y})\big]$.
}
\subsubsection{Función de Utilidad Esperada  de VON NEUMANN-MORGENSTERN}
En definitiva, si las preferencias del consumidor sobre las loterías son racionales, contínuas e independientes, entonces pueden ser representadas mediante una Función de Utilidad con la forma de Utilidad Esperada \footnote{%
    Por el Teorema de la Utilidad Esperada. La relación matemática entre una Función de Utilidad sin o con forma de utilidad esperada se puede ver en el Anexo, se trata de unas condiciones adicionales.
} 

\textbf{\textcolor{red}{[Definición: Matemáticas. Proposición 6.B.3]}}

Cuyas propiedades serán:

\begin{lnum}
    \item Preservará el orden de las preferencias;
    \item Es separable y aditiva; y,
    \item Es lineal en probabilidades.
\end{lnum}

\clearpage
\section{Loterías sobre dinero y actitud frente al riesgo}
Una de las aplicaciones más habituales del Teorema de la Utilidad Esperada es el análisis de las elecciones del consumidor en contextos de incertidumbre/riesgo donde el resultado son \textbf{cantidades monetarias}. 

\subsection{Definiciones y axiomas}
Para ello, partiremos de las siguientes premisas:

\begin{lnum}
    \item Consideraremos la cantidad de dinero como una variable continua;
    \item Describimos una \textit{lotería de dinero} como una función de distribución acumulada F\footnote{%
        Esto no excluye a \textit{priori} la posibilidad de que el codominio (recorrido) sea un conjunto discreto.
    }\footnote{%
        Recordar que las funciones de distribución, igual que las funciones de densidad, preservan la estructura lineal de las \textit{loterías}.
    };y,
    \item Un consumidor con preferencias racionales (completas y transitivas)
\end{lnum}
La aplicación del Teorema de la Utilidad Esperada a un conjunto de resultados continuo nos indica que, dados los supuestos base del Teorema, hay una asignación de valores de utilidad a valores no-decrecientes de dinero con la propiedad de que:

\textbf{[Proposición 6.C.1]}

En términos económicos podemos establecer una relación con lo expuesto aquí y los resultados obtenidos en el desarrollo de los temas [Tema 3.A.8 y 3.A.9]. La teoría de la demanda de consumo establece que dados los precios de los bienes, el equilibrio del consumidor es único (SCEDM o SCEDC) para cada nivel de renta \footnote{Tal y como se recoge en la Función Indirecta de Utilidad [Tema 3.A.8].}, luego cada nivel de renta monetaria tendrá asociada una solución óptima única de cesta/s de consumo en equilibrio\footnote{%
    A pesar de que se está analizando el comportamiento del consumidor, dado que este análisis parte de la existencia de un único bien o commodity, el mismo puede aplicarse al comportamiento de cualquier otro tipo de agente económico. Más concretamente, puede utilizarse para describir la conducta de la empresa o unidad de producción de bienes en situación de riesgo, con las debidas modificaciones (maximización del beneficio frente a maximización de utilidad, etc.).
}.

A pesar de que la axiomática general del Teorema de la Utilidad Esperada condiciona la representación (estructura) de la Utilidad Esperada. No establece ninguna restricción sobre la Función de Utilidad de Bernoulli. En gran parte, el poder analítico de la Utilidad Esperada proviene de la especificación de la Función de Utilidad de Bernoulli, de tal manera que capture atributos interesantes del comportamiento económico. Dicho esto, por razón de brevedad nosotros si estableceremos restricciones sobre la Función de Utilidad de Bernoulli, en concreto, diremos que es:
\begin{lnum}
    \item Creciente;
    \item Continua \footnote{%
        Estas dos condiciones son básicas y normales puesto que suponen caraterizar al agente como “Amante de la riqueza”. Nada raro.
    }
\end{lnum}

\imagenfit{imagen2.png}{[NO TITLE REQUIRED]}{[NO SOURCE REQUIRED]}

[1] - Para entender la representación en el símplex con triángulo rectángulo: Machina, M. J. (2018). Expected Utility Hypothesis. En Macmillan Publishers Ltd (Ed.), The New Palgrave Dictionary of Economics (p. 4204). Palgrave Macmillan UK. https://doi.org/10.1057/978-1-349-95189-5\_127]

Pasaremos ahora a uno de los resultados más interesantes de este planteamiento: \textit{la aversión al riesgo}; formulada a través de la Función de Utilidad de Bernoulli y sus mediciones\footnote{la aversión al riesgo}.

Para estudiar la aversión al riesgo del agente podemos acudir a sus dos caracterizadores más habituales, la concavidad o convexidad, por un lado, y el equivalente cierto, por otro, siendo ambos requisitos perfectamente equivalentes ,como veremos.

\subsubsection{Concavidad o convexidad de la Función de Utilidad de Bernoulli}
Utilizando la Desigualdad de Jensen\footnote{%
    La desigualdad de Jensen para funciones convexas relaciona el valor que asigna a una integral con la integral de esa misma función permutando, por así decirlo, la función y la integral. Fue probada por el matemático danés JOHAN JENSEN en 1906.
}, un agente será:
\begin{la}
    \item Averso al riesgo. Diremos que un individuo es averso al riesgo si su Función de Utilidad de Bernoulli $u(x)$ es cóncava, pues, \textit{in extremis}, valorará más la lotería degenerada \footnote{%
        Aquella que se encuentre en alguno de los vértices del Triplex, aunque sea menor [valor seguro $P(X=x) = 1$].
    }  que tenga como pago cierto su Utilidad Esperada (Expected Utility).
     \textbf{Es imprescinidible graficar esto}: $\int u(x)\,dF(x) \;<\; u\left(\int x\,dF(x)\right)$
    \item Neutral al riesgo. Por otro lado, diremos que un individuo es \textit{neutral al riesgo} si es indiferente entre escoger la lotería y $ \tilde{y}$ o la lotería degenerada que tenga como pago cierto su Utilidad Esperada. Por lo tanto, un agente neutral al riesgo  presentará una Función de Utilidad de Bernoulli lineal.
    \textbf{Es imprescinidible graficar esto}: $\int u(x)\,dF(x) \;=\; u\left(\int x\,dF(x)\right)$
    \item Amante del riesgo.  \textit{Senso contrario}, un individuo que prefiera recibir la lotería $ \tilde{y}$ a una lotería degenerada que tenga como pago cierto su esperanza será \textit{amante del riesgo}. Y, por tanto, su Función de Utilidad de Bernoulli será convexa.
    \textbf{Es imprescinidible graficar esto}: $\int u(x)\,dF(x) \;>\; u\left(\int x\,dF(x)\right)$
\end{la}
Un corolario directo de estas observaciones es el siguiente:

\textbf{Corolario}.- Para calificar la actitud frente al riesgo de cualquier individuo, simplemente basta comparar la Utilidad Esperada de una lotería (distrución de probabilidad) cualquiera con la Utilidad de la Esperanza de dicha lotería (distrución de probabilidad).

La intuición económica subyacente es de gran importancia: \textit{la aversión al riesgo es el sentimiento que el agente siente al valorar entre una pérdida y una ganancia}\footnote{%
    La intuición económica subyacente es de gran importancia: la aversión al riesgo es el sentimiento que el agente siente al valorar entre una pérdida y una ganancia
}. Concepto sobre el que luego volveremos.

\subsubsection{Equivalente cierto y prima de riesgo}.
Como segundo caracterizador veremos la prima de riesgo. Para ello, es necesario definir previamente el concepto de equivalente cierto de una lotería.
[\textbf{Definimos matemáticamente. Definición 6.C.2}]

En palabras, el equivalente cierto es la cantidad de dinero que debe percibir un individuo para que le sea indiferente apostar (elegir una lotería) o escoger aquella cuyo equivalente cierto sea la EU.

(\textbf{Todo lo que sigue es consecuencia directa de las definiciones previas, de hecho, hay sobrenotación (falta especificar la utilidad de la esperanza de $E[x]$): revisar el tema de CECO})

Por tanto, complementando las definiciones anteriores podemos decir que:
\begin{la}
    \item El individuo será averso al riesgo  cuando el equivalente cierto sea menor al pago monetario esperado de la lotería.
    $\xi < E[x], \quad \forall F(x)$
    \item b)	El individuo será neutral al riesgo cuando el equivalente cierto sea igual al pago monetario esperado de la lotería. 
   $\xi < E[x], \quad \forall F(x)$
    \item El individuo será amante del riesgo cuando el equivalente cierto sea mayor al pago monetario esperado de la lotería.  
    $\xi < E[x], \quad \forall F(x)$
\end{la}

Por otro lado, definiremos prima de riesgo como la diferencia entre el valor esperado de la lotería $E[x]$ y el equivalente cierto $\xi$, y nos da una medida monetaria de la cantidad que está dispuesto a pagar un agente para evitar un riesgo\footnote{%
    O, equivalentemente y con aplicación en la teoría de la selección de carteras, la cantidad que está dispuesto a aceptar para tomar un riesgo.
}\note{Borrador del Manual de Macro de Jesús Fernández Villaverde: No lo usé mucho pero tiene algún detalle interesante más allá del álgebra de los modelos. Por ejemplo, cuando plantea el modelo de crecimiento te explica muy bien las funciones CRRA (aversión relativa constante al riesgo). De hecho, en este enlace en el apartado 3.1. te explica algo muy interesante sobre el comportamiento precautorio y la aversión al riesgo: https://www.sas.upenn.edu/\~jesusfv/Uncertainty_Shocks_Business_Cycle.pdf (págs. 9, 10 y 11) Ver esto para entender lo de LELAND (1968) y SANDMO (1970) – Útil para [Tema 3.A.33] de demanda de consumo en la teoría del paseo aleatorio de HALL (1978). En este artículo buscan explicar mecanismos para explicar cómo shocks de incertidumbre pueden causar ciclos económicos. En el apartado 3.1. destacan el papel del comportamiento precautorio, de forma que si las preferencias de los agentes del modelo muestran comportamiento precautorio, un shock de incertidumbre puede causar ciclos económicos. Cuando la función de utilidad de un agente económico es cóncava, el agente es averso al riesgo y no le gusta la incertidumbre. En este caso, la utilidad esperada del consumo es más baja que la utilidad del consumo esperado (tal y como hemos visto en el Tema 3.A.10). La aversión al riesgo es condición necesaria pero no suficiente para el comportamiento precautorio. Este quiere decir, que dada la existencia de riesgo, el agente cambia su comportamiento como consecuencia del mismo. Para que exista comportamiento precautorio, la tercera derivada ha de ser mayor que cero, es decir, la convexidad de la utilidad marginal es condición suficiente para comportamiento precautorio [también conocido como prudencia]. Esto se da una función de aversión constante al riesgo (CRRA). Sin embargo, en una función de utilidad cuadrática a pesar de caracterizar a un averso al riesgo (ya que $u''(x)<0$), el agente no muestra comportamiento precautorio, es decir, no cambia su comportamiento como respuesta al riesgo.
¿Cuáles son las consecuencias a nivel agregado de estos cambios en el comportamiento precautorio? En un modelo macroeconómico base, con una función de utilidad CRRA, el comportamiento precautorio aparece típicamente como ahorro precautorio y por tanto un aumento en la incertidumbre conduce a una mayor demanda de ahorros, reduciendo el tipo de interés real de equilibrio. Esta reducción del tipo de interés induce un mayor consumo presente, por lo que se produce un efecto contrario a los efectos contractivos del ahorro precautorio. Por tanto, el efecto final queda indeterminado. Esto será relevante en la teoría de la demanda de consumo del paseo aleatorio de HALL:
\imagenfit{img1_p46.png}{[NO TITLE REQUIRED]}{[NO SOURCE REQUIRED]}}.

$\rho = E[x] - \xi$

\subsubsection{Coeficientes de aversión al riesgo}
Una vez delimitadas las tres grandes categorías de agentes por su actitud frente al riesgo, se podría tratar de refinar un poco más el análisis buscando una medida que cuantifique esta actitud, permitiendo cuantificar la aversión al riesgo de un agente. Es por esta razón que surgen los coeficientes de aversión al riesgo:
\begin{lnum}
    \item En primer lugar, analizaremos el coeficiente de aversión absoluta al riesgo de ARROW-PRATT. Dada una Función de Utilidad de Bernoulli, dos veces diferenciable\footnote{%
        Esta medida recoge la idea de que el grado de aversión al riesgo puede medirse por el grado de curvatura de la Función de Utilidad de Bernoulli, pues cuanta más curvatura más nos alejamos del caso de neutralidad al riesgo (curvatura nula).
    }:
    \item[--] \hspace{-2em} $r\alpha = -\,u''(x)\,u'(x)$
    \item[--] \hspace{-2em} A fin de hacer la \textit{curvatura de la función} (numerador) invariante ante transformaciones lineales crecientes de la función, es necesario dividirla por la primera derivada, tal y como vemos, añadiendo el cambio de signo sirve para obtener una medida creciente con dicho grado de aversión..
    \item[--] \hspace{-2em} El coeficiente de aversión absoluta al riesgo de ARROW-PRATT permite realizar comparaciones entre individuos y entre distintos niveles de renta para un mismo individuo (pues nada implica que el coeficientesea constante, pudiendo variar atendiendo al nivel de renta).
    \item[--] \hspace{-2em} Esta particularidad lleva a refinar un poco más el coeficiente de aversión, pues es sensato pensar que los agentes de mayor riqueza estarán dispuestos a aceptar mayores riesgos en términos absolutos que los de escasa riqueza, es decir, que los agentes exhiban aversión al riesgo decreciente en renta (decreasing absolute risk aversion o DARA).
    \item Para depurar el coeficiente de los efectos de diferentes niveles de renta puede usarse el denominado coeficiente de aversión relativa al riesgo, que mide dicha aversión en términos de ganancias o pérdidas porcentuales de riqueza y se define como:
    \item[--] \hspace{-2em} $rr = r\alpha \cdot x = -\,u''(x)\,u'(x)\cdot x$
\end{lnum}
Como puede observarse, el coeficiente relativo es una elasticidad de la utilidad marginal\footnote{%
    $ r_r = r_\alpha \cdot x = -\frac{u''(x)}{u'(x)} \cdot x = -\frac{\partial u'(x)}{\partial x} \cdot \frac{x}{u'(x)} = \varepsilon_{u'(x),x}$
}  y como tal no está afectada por las unidades de medida de la riqueza. Respecto a la misma, en muchas aplicaciones es habitual suponer que los agentes exhiben aversión relativa al riesgo no creciente \footnote{%
    Esto es, se presupone únicamente su monotonía negativa, pero no necesariamente estricta. Puede haber sujetos conservadores muy ricos.
} (\textit{non-increasing relative risk aversion})\note{Esto no es necesario incluirlo ahora, reulta muy complicado sin la relación de propiedades y condiciones que corresponden a otro tema. Si se quiere incluir, añadir de forma anecdótica]
De hecho, una categoría de funciones muy utilizada en macroeconomía (por ser muy habitual y muy conveniente) es la función de aversión relativa al riesgo constante (Constant Relative Risk Aversion, CRRA) de la siguiente forma\footnote{%
    En macroeconomía se trabaja mucho con una función isoelástica (Constant Relative Risk Aversion - CRRA). Esto es una propiedad atractiva, ya que a lo largo de los dos últimos siglos, los niveles de prima de riesgo asociados al comercio de activos financieros en EE.UU. y en Reino Unido han fluctuado alrededor de una media considerablemente constante a pesar del crecimiento de la renta per cápita y del consumo de ambos países. Además, a nivel analítico permite obtener primas de riesgo estacionarias cuando el consumo crece a una tasa constante. La función isoelástica también tiene como implicación (con adicionales supuestos) que la selección de cartera de un hogar (i.e. qué parte de su riqueza invertir en acciones o en bonos) no se ve alterada por cambios en la riqueza del hogar.
}:
\[ u(x) = \begin{cases} \dfrac{x^{1-\theta}-1}{1-\theta}, & \text{si } \theta \geq 0 \text{ y } \theta \neq 1, \\
\ln(x), & \text{si } \theta = 1. \end{cases} \]
El coeficiente de aversión relativa al riesgo ARROW-PRATT de esta función es constante:\\
ECUACION PENDIENTE\\
Además, esta forma funcional es muy utilizada en macroeconomía por presentar elasticidad de sustitución intertemporal constante. Para este tipo de funciones, se demuestra que ESI≡σ=1θ⁄, siendo θ el parámetro de la función de utilidad que, a su vez, representa la elasticidad de la utilidad marginal [ver Tema 3.A.29]\footnote{%
    2 La elasticidad de sustitución intertemporal quedaría definida como [ver Tema 3.A.29]: [FÓRMULA]
De este modo, con funciones de utilidad separables en el tiempo, la inversa de la elasticidad de sustitución y el coeficiente de aversión relativa al riesgo son idénticos. Por lo tanto, la familia de funciones de utilidad con aversión relativa al riesgo constante también consiste de aquellas funciones con elasticidad de sustitución intertemporal constante.
}.}.
Todos estos resultados se pueden ver gráficamente, en el caso de dos resultados posibles:
\imagenfit{img3_p16.png}{[NO TITLE REQUIRED]}{[NO SOURCE REQUIRED]}
\subsection{Aplicaciones: la contratación de seguros}
[Mencionar que tiene varias extensiones pero solo cantar una por cuestiones de tiempo. Quizás sea más simple y más intuitivo la demanda de seguros, que también es un paso previo al desarrollo del modelo de ROTHSCHILD y STIGLITZ del Tema 3.A.13]

Vista la herramienta básica, es necesario estudiar la aplicación de la utilidad esperada a problemas económicos concretos. En general, las decisiones más importantes del consumidor en situación de riesgo son las decisiones financieras, pues son aquellas en las que el riesgo es un aspecto ineludible.

Por decisiones financieras se entienden aquellas que determinan la demanda por el consumidor de diferentes instrumentos financieros (en este caso, activos financieros) que en última instancia son siempre “promesas” de pagos (típicamente inciertos), que se diferencian entre sí básicamente por el evento (más o menos incierto) a que se condiciona ese pago prometido.

Así, la teoría de la utilidad esperada encuentra aplicaciones en temas como:

\begin{lnum}
    \item Las decisiones de cartera; o,
    \item La decisión de aseguramiento, que es en la que nos centraremos en esta exposición.
\end{lnum}
\subsubsection{Modelo}
\textbf{Idea}

Este modelo, propuesto inicialmente por ROTSCHILD y STIGLITZ (1976) consiste en la resolución de elección o no de un asegramiento, cuyo fin será reducir su exposición al riesgo (expresado mediante la posibilidad de perder una cantidad determinada de su renta original).

El agente se encuentra ante un continuo de resultos posibles (tantos como niveles de renta), pero simplificaremos el modelo a dos posibles resultados con probabilidad no nula. 

\textbf{[DEFINIR DOS RESULTADOS]}

\begin{lnum}
    \item Resultado 1 sin contingencia (representado en abcisas); y,
    \item Resultado 2 con contingencia (p.e. el robo o pérdida de una parte de su riqueza; en ordenadas).
\end{lnum}
Dicha pérdida contingente la expresaremos como \text{P}, de modo que su renta sería \textit{W-P}.

\textbf{[DEFINIR PROBABILIDADES ASIGNADAS Y SU DISTRIBUCIÓN]}

Esto implica que el agente tiene su punto de partida en un punto como A(W, W-P), en que el valor de abscisas es mayor que el de ordenadas \footnote{%
    Nótese que el lugar geométrico que representa la ausencia de riesgo es la bisectriz del cuadrante, conocido como \textbf{linea de certeza}, pues sobre ella la renta es la misma independientemente de la realización del resultado
}.

Las preferencias del agente vendrán representadas por una función de utilidad esperada v.N-M

\textbf{[DEFINIR LAS PREFERENCIAS DEL AGENTE]}

Y consideraremos que su Función de Utilidad de Bernoulli será continua, creciente y cóncava, indicando esto último que el agente es averso al riesgo. Por tanto, sus curvas de indiferencia serán:
\begin{lnum}
    \item Continuas;
    \item Decrecientes;
    \item Representando más utilidad esperada conforme se alejan de los ejes;
    \item Convexas respecto al origen de coordenadas\footnote{%
        Obsérvese que esta utilidad von Neumann-Morgenstern se asemeja mucho a una función Cobb-Douglas linealmente homogénea sometida a la habitual transformación logarítmica que simplifica su resolución.
    }; y
    \item Simétricas respecto a la linea de certeza\footnote{%
        Por ser esta Función de Utilidad de Bernoulli estado-independiente (\textit{state-independant}).
    } 
\end{lnum}
\textbf{[DIBUJAR GRÁFICO DEL PROBLEMA]}

Con la ayuda de este gráfico podemos también aprovechar para definir la \textbf{Relación Marginal de Sustitución Estocástica (RSME)}, o tasa a la que un individuo está dispuesto a sacrificar renta en un resultado a cambio de icrementarlo en el otro resultado\footnote{%
    La RSME viene definida por el cociente de la derivada parcial de la Función de Utilidad de Bernoulli respecto al \textit{Resultado 1}, entre la derivada parcial respecto al \textit{Resultado 2}. La RSME no es más que la RMS de certidumbre multplicada por el cociente de probabilidades.
}. Además, se oberva que todas las curvas de indiferencia, al cortar a la línea de certeza, poseen la misma pendiente (cociente entre la probabilidad del Resultado 1 entre la probabilidad del Resultado 2). Por lo tanto, por el punto de dotación inicial pasará una curva de indiferencia que dará el nivel de utilidad esperado en dicho punto.

Desde esta situación inicial, se le ofrece aseguramiento de su riqueza a cambio de una prima, \textit{r}, proporcional a la cantidad asegurada, \textit{S} (cantidad no determinada\footnote{%
    Pudiendo ser superior, inferior o igual a la cantidad \textit{P} que esté en juego.
}\footnote{%
    Nótese que la variable de decisión para el agente es la cuantía asegurada \textit{S}, lo que en el presente modelo es equivalente a elegir una lotería. En efecto, al elegir \textit{S} se eligen indirectamente los niveles de renta que se disfrutarán en último término en cada resultado, con lo que la lotería cambiará al cambiar los niveles de renta a los que se asigna una probabilidad dicha probabilidad a sucesos (niveles de renta) diferentes.
}). De este modo, el individuo obtendrá una utilidad:

\textbf{[DEFINIR LAS RESTRICCIONES DEL PROBLEMA DE OPTIMIZACIÓN Y OBTENER LINEA DE ASEGURAMIENTO (IGUALAR LAS RESTRICCIONES RESPECTO A S, DEPEJAR \textit{RESULTADO 2}). DIBUJAR SOBRE EL GRÁFICO ANTERIOR LA LINEA DE ASEGURAMIENTO]}

De la recta de aseguramiento descubrimos que el precio relativo o relación de intercambio que el agente debe soportar para trasladar renta entre distintos resultados es\footnote{%
    La pendiente negativa de esta cuva está indicando que el agente, para poder aumentar su renta en caso de contingencia, reduciendo su eventual pérdida ($dx_2$), debe renunciar a renta también en el caso de no-contingencia ($-dx_1$).
}:

\textbf{[INTRODUCIR ECUACIÓN]}

El equilibrio del agente se obtiene resolviendo el problema de optimización siguiente:

\textbf{[PLANTEARLO]}

En términos formales, el equilibrio debe darse con las siguientes CPO’s\footnote{%
    Se pueden dar también soluciones de esquina, en que el equilibrio se alcanza sin darse esta igualdad ya sea porque el precio que el agente debe soportar (primas) es demasiado alto y prefiere quedarse en su situación de partida (en cuyo caso, RMSE > precio relativo), ya sea porque la prima le resulta demasiado barata hasta el punto de aceptar un riesgo más allá de la situación de seguro pleno (RMSE < precio relativo).
} de forma que el agente iguala la tasa a la que está dispuesto a intercambiar renta entre los distintos resultados (RSME) con la tasa objetiva o precios relativos de mercado. 

\textbf{[PLANTEAR SISTEMA]}

En términos gráficos, la CPO anterior es una condición de tangencia entre la curva de indiferencia más elevada posible y la \textit{recta de aseguramiento}, obtenida mediante la igualación de las pendientes de ambas.
\subsubsection{Implicaciones: el tipo de prima}
Aunque la resolución completa del ejercicio acabará con una demanda de aseguramiento \textit{\*S}, que analíticamente dependerá de un conjunto de factores y parámetros\footnote{%
    Se pueden dar también soluciones de esquina, en que el equilibrio se alcanza sin darse esta igualdad ya sea porque el precio que el agente debe soportar (primas) es demasiado alto y prefiere quedarse en su situación de partida (en cuyo caso, RMSE > precio relativo), ya sea porque la prima le resulta demasiado barata hasta el punto de aceptar un riesgo más allá de la situación de seguro pleno (RMSE < precio relativo).
  Entre los que están unos relativos a la Función de Distribución de Bernoulli y otros que afectan a la forma concreta que adopta una lotería o función de distribución (y por ende su media, dispersión, asimetría, \textit{kurtois}, etc.):
  \begin{lnum}
      \item El grado de aversión al riesgo del agente, sintetizado a través del coeficiente de aversión al riesgo, afectará positivamente a la demanda de aseguramiento;
      \item La probabilidad de contingencia, que implicaría (en la presente especificación) una reducción de la RMSE y un aplanamiento de las curvas de indiferencia, lo que provoca que una prima dada sea más favorable, incentivando una mayor demanda de aseguramiento;
      \item La renta del agente, \textit{W}, aunque su efecto es \textit{a priori} indeterminado ya que se debería concoer el tipo de aversión al riesgo que presenta el agente, se demuestra (GRAVELLE y REES, 2004) que la demanda aumenta o disminuye con la renta según el agente exhiba aversión absoluta;
      \item La magnitud de la pérdida, \textit{P}, también afectará positivamente a \textit{S}, \textit{ceteris paribus}. 
      \item La prima o precio de aseguramiento, r, el cual es razonable pensar que afectará a la demanda de aseguramiento. 
  \end{lnum}
}.

\textbf{[REPRESENTAR LA FUNCIÓN S* EN FUNCIÓN DE LOS DIFERENTES PARÁMETROS]}

El parámetro más interesante es el precio o prima unitaria r, que puede clasificarse en tres niveles diferentes permitiendo predecir el nivel de aseguramiento de cualquier agente. 
\begin{la}
    \item \textit{Prima actuarialmente favorable (S > P).} El individuo se sobreaseguraría, es decir, la recta de aseguramiento tiene mayor pendiente que la curva de indiferencia en el punto de corte con la linea de certeza, con lo que el equilibrio se situaría por encima (a la izquierda) de la linea de certeza.
    \item \textit{Prima actuarialmente justa (S = P).} El individuo se aseguraría completamente, es decir, el quilibrio se sitúa sobre su linea de certeza (el agente iguala entre resultados). Supuesto de competencia perfecta en el mercado de seguros.
    \item \textit{Prima actuarialmente desfavorable (S < P).}\footnote{%
        Caso extremo (con solución de esquina), no contrará ni siquiera el seguro y permanecerá en su punto de dotación inicial A.
    } La recta de aseguramiento tiene menor pendiente que la curva de indiferencia en el punto de corte con la linea de certeza. 
\end{la}

Dados los supuestos iniciales de nuestro modelo, se da la situación E, en la que el individuo se asegura completamente maximizando su utilidad y las empresas obtienen beneficio nulo (competencia perfecta):

\imagenfit{img7_P20.png}{[NO TITLE REQUIRED]}{[NO SOURCE REQUIRED]}

Sin embargo, los resultados variarían si modificamos, por ejemplo, el tipo de aversión al riesgo de nuestro agente:

 \imagenfit{img1_P21.png}{[NO TITLE REQUIRED]}{[NO SOURCE REQUIRED]}
 
 \imagenfit{img2_P21.png}{[NO TITLE REQUIRED]}{[NO SOURCE REQUIRED]}
 
\subsubsection{Valoración}
Este modelo permite tener una visión general con los fundamentos de la demanda de aseguramiento\footnote{%
    Es interesante conocer, a efectos de otros temas, en qué consiste un reparto óptimo del riesgo (entre agentes con diferentes preferencias frente al riesgo, por ejemplo), así como ciertos fenómenos que se producen en la industria aseguradora cuando se cubre a una colectividad de agentes (\textit{risk pooling} o agrupación en un colectivo de individuos de sus riesgos individuales, por ejemplo, una mutua de seguros) y cuando existe una colectividad de aseguradores (\textit{risk spreading} o distribución del riesgo entre numerosos aseguradores aversos al riesgo que conjuntamente se comportarán como un agente neutral al riesgo, bajo ciertas hipótesis, según el teorema de ARROW y LIND). Aunque caen fuera del ámbito de la elección del consumidor, se remite para su explicación al manual de GRAVELLE y REES (2004).
} para conocer sus determinantes, sus posibles equilibrios o resultados esperados, así como el análisis de estática comparativa.

No obstante, la simplicidad de sus supuestos lo hacen susceptible de importantes extensiones o refinamientos que modifican sus supuestos de partida, entre las que destaca por su relevancia la existencia de información asimétrica, tanto ex-ante que origina problemas de selección adversa, como ex-post con problemas de riesgo moral [ver Tema 3.A.13].
\subsection{Decisión de cartera: demanda de bonos y acciones}
[Lo mismo pero al revés, cuánto es necesario para apartar a un individuo averso al riesgo de su línea de certeza. No cantar. Simplemente mencionar.]
\textbf{[Extensiones en Víctor Gutiérrez-Marcos]}
\subsection{Dominancia estocástica}
\textbf{[Todo lo expuesto anteriormente restringio a las loterías con carteras/dinero]}
Pero ¿existe una forma objetiva (sin necesidad de utilizar funciones de utilidad subjetivas) de comparar alternativas de riesgo? Es decir, el propósito aquí sería buscar caracterizadores de esas loterías o distribuciones de probabilidad, lo suficientemente sencillos o manejables, que permitan clasificarlas de manera objetiva.
\subsubsection{Dominancia estocástica de primer orden}
En primer lugar, la pregunta sería: ¿existen loterías que son preferidas a otras para cualquier agente independientemente de sus preferencias, esto es, independientemente de $u(x)$? Responder a esta pregunta lleva a categorizar las distintas loterías a través del concepto de dominancia estocástica de primer orden\footnote{.} 28. Se dice que una distribución $F(x)$ domina estocásticamente en primer orden a otra $G(x)$ si es preferida por cualquier agente amante de la riqueza independiente de su actitud frente al riesgo, es decir, si:\\
$\int u(x)dF(x)\geq\int u(x)dG(x)\ \forall\ u(x)$  tal que $u'(x)> 0$\\
De este modo, \textbf{todos los amantes de la riqueza $(u'(x)> 0)$ prefieren siempre $F(x)$}.

[28 - Formalmente se tiene que la lotería $F$ domina a la lotería $G$ estocásticamente de primer orden si para cada $x \in \mathbb{R}$, se verifica que $F(x) \leq G(x)$. Además, si $F$ domina a $G$ estocásticamente de primer orden, entonces cualquier agente con preferencias crecientes prefiere $F$ a $G$.

Ejemplo gráfico: En el siguiente gráfico, la lotería $F$ domina estocásticamente de primer orden a la lotería $G$.

\imagenfit{img1_P23.png}{[NO TITLE REQUIRED]}{[NO SOURCE REQUIRED]}

\subsubsection{Dominancia estocástica de segundo orden}
La dominancia estocástica de segundo orden\footnote{.} 29 se utiliza para comparar el “riesgo” de 2 loterías.

Para concentrarnos sólo en ese aspecto de las loterías supondremos que todas tienen el mismo valor esperado (ya que un agente puede aceptar un riesgo mayor si el pago esperado es suficientemente grande). \textit{De este modo, dadas 2 loterías con igual esperanza pero distintas distribuciones, no podremos afirmar que una lotería domine estocásticamente de primer orden a la otra, pero podremos estudiar si cualquier persona aversa al riesgo elegiría una lotería sobre la otra. Esto es lo que se conoce como dominancia estocástica de segundo orden} \textbf{[Creo que no es verdadero]}.

En general, se dice que una distribución $F(x)$ domina estocásticamente en segundo orden a otra $G(x)$ si es preferida por cualquier agente averso al riesgo, es decir, si:
$\int u(x)\, dF(x) \;\geq\; \int u(x)\, dG(x) \quad \forall\, u(x) \text{ tal que } u'(x) > 0 \;\text{ y }\; u''(x) < 0$
De este modo, ya no todos los amantes de la riqueza $(u'(x)>0)$ prefieren siempre $F(x)$, pero aún es posible afirmar que todos los aversos al riesgo $(u"(x)<0)$ prefieren la lotería $F(x)$ a la $G(x)$. De este modo, ya no todos los amantes de la riqueza $(u'(x)>0)$ prefieren siempre $F(x)$, pero aún es posible afirmar que todos los aversos al riesgo $(u"(x)<0)$ prefieren la lotería $F(x)$ a la $G(x)$.

Desgraciadamente, como se ha comprobado, las propiedades de dominancia estocástica de primer y de segundo orden, aun basándose únicamente en las nociones de esperanza y dispersión, tienden a ser poco prácticas por demasiado exigentes y, en el mejor de los casos, solo permiten ordenar las diferentes loterías disponibles\footnote{.}30 pero no cuantificar cuánto mejor es una lotería que otra.

[30 - Se basan en una función de utilidad esperada que es de naturaleza ordinal.]

Es posiblemente por esta razón que se acuda a modelos de decisión simplificados basados en el binomio media-varianza como versión reducida\footnote{.}31 del modelo de utilidad esperada von Neumann-Morgenstern, representando la media un “bien” y la varianza un “mal” (se asume \textbf{aversión al riesgo}).

[31 - Nótese que el modelo media-varianza descansa únicamente sobre un momento de primer orden (media) y otro de segundo orden (varianza), omitiendo momentos de orden superior que como ya se ha dicho pueden afectar a las decisiones de los agentes al reflejar la asimetría (momentos de tercer orden) y la curtosis (momento de cuarto orden de las distribuciones de probabilidad.]

Como apuntan BICKCHANDADI, HIRSHLEIFER y RILEY (2013), el paso desde una distribución de probabilidad completamente definida (en términos de las consecuencias para todos y cada uno de los estados) a un mero resumen estadístico de la distribución (que sólo tiene en cuenta dos parámetros: media y varianza) implica necesariamente una pérdida de información. La cuestión que se plantea es bajo que supuestos dicha simplificación es válida.

Con carácter general, una posible justificación es el Teorema Central del Límite, de la teoría de la probabilidad que, dicho muy simplificadamente, predice que la suma de cualquier número elevado de N variables aleatorias converge hacia la distribución normal, una distribución que precisamente posee la propiedad de quedar plenamente descrita con solo dos parámetros: media y varianza.

El campo más propicio de los modelos de media-varianza es por tanto las situaciones en que se da dicha suma de múltiples variables aleatorias, siendo el ejemplo más paradigmático la teoría de selección de carteras desarrollada, entre otros, por MARKOWITZ, TOBIN y SHARPE [ver tema 3.B.23].

IMG 4 P25 

[32 - https://www.youtube.com/watch?v=pNmX7g-FRrw]2

\subsection{Extensiones y primeras críticas (a la TUE)}
\subsubsection{Extensiones}
El enfoque de utilidad esperada VON NEUMANN-MORGENSTERN es el más extendido, pero también es susceptible de ciertas extensiones, entre las que destacan:
\begin{lnum}
    \item La introducción de \textbf{utilidad estado-dependiente} (\textit{state-dependant utility}). En el modelo básico estudiado se asume que la utilidad de Bernoulli es la misma independientemente del resultado (estado-independiente). Sin embargo, no sería rao pensar que dicha función de utilidad varía también con el resultado. Es decir, sería sensato pensar que la causa concreta que origina estos resultados (llamados estados de naturaleza) también afecta a la utilidad, que pasaría a definirse para cada-estado en particular. El ejemplo más típico es el caso en que los diferentes estados de naturaleza se refieren a la salud, pues no se disfrutaría igual un mismo pago con buena salud que con la salud mermada. Sin entrar en especificaciones de modelos, podemos simplementer comentar que existe una extensión del Teorema de la Utilidad Esperada (\textit{“Teorema de la Utilidad Esperada Ampliada o Extendida”}) que demuestra la existencia de ésta, en términos análogos a la previamente propuesta. 
    \item El \textbf{enfoque de preferencia por el estado} (\textit{state-preference}) planteada por ARROW (1953) y DEBREU (1959) permite reducir el problema de elección bajo riesgo a un problema de elección convencional (bajo certidumbre) sobre el espacio de resultados posibles, mediante un cambio adecuado de la estructura de los bienes. De esta manera, las preferencias del consumidor se formarían sobre \textbf{cestas estado-contingentes}, i.e. no solo por sus propiedades físicas (como en el caso de certidumbre) sino también por el momento y el estado de la naturaleza futuro en que se consuman, siendo los precios de éstos conocidos como \textbf{precios Arrow-Debreu}. Podríamos decir, en linea con lo expuesto sobre variedades en el [Tema 3.A.8] que para cada bien físico existen tantas variedades como momomentos del tiempo y estados de naturaleza en que se consume, debiendo existir mercados diferentes para cada uno de ellos: \textit{“No se valora lo mismo un paragüas cuando llueve que un paragüas cuando hace sol”}. Este enfoque permite incorporar fácilmente el riesgo en modelos de equilibrio general, pero la carencia de realismo sobre sus supuestos limita su aplicación práctica. 
\end{lnum}
Hasta aquí hemos desarrollado la Teoría de la Utilidad Esperada y hemos visto cómo modeliza la teoría económica el comportamiento de los individuos cuando estos se enfrentan a situaciones de riesgo o incertidumbre.

La Teoría de la Utilidad Esperada representa la herramienta de análisis fundamental para el estudio de situaciones de riesgo. Su conveniencia y operatividad son innegables.
\subsubsection{No obstante, la Teoría de la Utilidad Esperada no está exenta de inconvenientes:}
No obstante, \textbf{la Teoría de la Utilidad Esperada no está exenta de inconvenientes}:

Entre las primeras críticas destaca el análisis de FRIEDMAN y SAVAGE (1948), que apuntaron la aparente paradoja que suponía que un mismo agente decidiera adquirir seguros contra contingencias (lo que revela aversión al riesgo), y a la vez participar en juegos de azar que son abiertamente desfavorables en términos estadísticos (lo que revela amor por el riesgo).

La solución propuesta por estos autores\footnote{%
    Que deja entrever su \textit{amor o afinidad} por este planteamiento, pues no rebaten sus bases sino que simplemente ofrecen un \textit{“ajuste”} metodológico para acomodar esta paradoja.
} fue que la Función de Utilidad de Bernoulli poseyera varios tramos de concavidad y convexidad con dos puntos de inflexión (cóncavo – convexo – cóncavo), conjeturando que la aversión al riesgo se da en los niveles más altos y más bajos de renta, pero que sin embargo en los niveles intermedios se da amor por el riesgo. Desgraciadamente, se trataba de una solución incompleta y con contradicciones empíricas como que los individuos más ricos nunca aceptarían apuestas justas. \textit{La existencia de casinos en Montecarlo invalida esta hipótesis}.

\textbf{[Mostrar planteamiento de Friedman y Savage. Graficar función de Bernoulli con diferentes tramos]}

Posteriormente, HARRY MARKOWITZ\footnote{%
    HARRY MARKOWITZ estuvo bajo la supervisión doctoral de FRIEDMAN y MARSCHAK fue galardonado con el Premio Nobel de Economía en 1990 junto con MERTON MILLER y WILLIAM SHARPE «Por sus trabajos pioneros para establecer la teoría de la economía financiera»..
}, criticó la Función de Utilidad FRIEDMAN-SAVAGE y argumentó en su artículo “The Utility of Wealth” (1952), que en la concavidad final de su función, se asume que los individuos con las rentas más altas nunca apostarán\footnote{Los casinos de Monte Carlo no tendrían entonces razón de ser.}. Para solucionar esta anomalía, MARKOWITZ amplía la crítica de FRIEDMAN y SAVAGE en dos aspectos.

La primera aportación consistía en añadir un tercer punto de inflexión. Así según HARRY MARKOWITZ, los individuos:

\imagenfit{img1_P27.png}{[NO TITLE REQUIRED]}{[NO SOURCE REQUIRED]}

\begin{lnum}
    \item Son aversos:
    \begin{enumerate} 
        \item a riesgos pequeños negativos\footnote{%
            El individuo no quiere aceptar un juego que pueda reducir un poco su riqueza (rechaza apuestas pequeñas desfavorables). La utilidad es cóncava alrededor del punto de referencia. Comportamiento prudente.
        } y
        \item a riesgos grandes positivos\footnote{%
            El individuo no se siente atraído por apuestas con grandes ganancias pero alta incertidumbre (quizá porque parecen irreales o demasiado lejanas). En niveles muy altos de ganancia, la función vuelve a ser cóncava, \textit{“saturación del placer marginal”}.
        }; y
    \end{enumerate}
    \item Son amantes:
        \begin{enumerate} 
        \item a riesgos grandes negativos\footnote{%
            El individuo acepta pequeñas apuestas favorables (por ejemplo, jugar una lotería barata o apostar algo trivial). La utilidad es convexa junto a la derecha del punto de referencia, se trata de un comportamiento \textit{buscador de riesgo en ganancias pequeñas}.
        }; y
        \item 	a riesgos pequeños positivos\footnote{%
            El individuo puede asumir riesgos desesperados para evitar pérdidas muy grandes o situaciones extremas (catástrofes financieras). En el tramo más a la izquierda, la función vuelve a ser convexa. El individuo busca el riesgo como último recurso.
        }.
    \end{enumerate}
\end{lnum}
La segunda aportación (y la más importante) de MARKOWITZ fue sugerir \textbf{que las loterías no se evalúan según el nivel de renta final, sino en función de los “cambios de renta” tomando como referencia su renta presente} [Esta idea, fue recogida posteriormente por la Prospect Theory de KAHNEMAN y TVERSKY\footnote{.}].

\clearpage
\section{Probabilidades subjetivas e incertidumbre}
Continuamos nuestro análisis preguntandonos: ¿Son siempre los riesgos conocidos? De hecho, rara veznos enfrentamos a situaciones de riesgo como las descritas, es decir, reducibles a distribuciones de probabilidad objetivas perfectamente definidas. En general, los agentes se enfrentan a situaciones cuya aleatoriedad no puede expresarse de manera cuantitativa y cuando el término “probabilidad” aparece, se refiere generalmente a imprecisas estimaciones subjetivas.

Por tanto, continuando con la delimitación knightiana, situaciones de incertidumbre serán todas aquellas en las que la aleatoriedad a la que se enfrenta el agente económico no se presenta en forma de probabilidades objetivas (ni especificadas exógenamente ni calculables científicamente).

Un primer intento de formalizar la incertidumbre en el sentido de KNIGHT es suponer que el agente se comporta como si estableciera probabilidades subjetivas o creencias, denominado enfoque de la probabilidad subjetiva, por lo que en esencia se puede utilizar el mismo aparato formal de la Teoría de la Utilidad Esperada.
\subsection{Orígenes de la Teoría de la probabilidad subjetiva}
Apenas unos años después de que VON NEUMANN y MORGENSTERN publicaran su modelo, SAVAGE (1954) aportó a la teoría de la utilidad esperada el análisis con probabilidades subjetivas, que generalizaba el modelo de von Neumann-Morgenstern para contextos en los que los agentes no conocen las distribuciones objetivas de probabilidad de los diferentes resultados posibles. Posteriormente se han desarrollado enfoques similares como el de ANSCOMBE y AUMANN (1963) y el de WAKKER (1989)\footnote{%
    Estos modelos difieren en la especificación de los conjuntos de elección y las correspondientes estructuras de preferencia, pero su objetivo es el mismo, la determinación simultánea de: \begin{lnum}
    \item Una función de utilidad que cuantifique los gustos del decisor; y
    \item Una medida de probabilidad que cuantifique sus creencias..
\end{lnum}
}, aunque se analizará brevemente el primero pues ser su primera aproximación.
\subsection{Modelo de SAVAGE}
\subsubsection{Planteamiento del Modelo}
El modelo de SAVAGE, se vale tanto de la axiomática básica de la TUE como del enfoque estado \textit{dependendiente} (de la TUE), y establece una serie de condiciones que debe cumplir la relación de preferencia entre actos (funciones que relacionan el conjunto de estados de la naturaleza con el de sus posibles consecuencias) para que el agente se comporte como si estuviese estimando implícitamente una probabilidad para cada estado de la naturaleza. Es decir, para que existan probabilidades subjetivas\footnote{%
    Karni, E. (2018). Savage’s Subjective Expected Utility Model. En Macmillan Publishers Ltd (Ed.), The New Palgrave Dictionary of Economics (pp. 11938-11942). Palgrave Macmillan UK. https://doi.org/10.1057/978-1-349-95189-5_2467 Fonseca, G. (2021) “Subjective Expected Utility”. The History of Economic Thought Website. https://www.hetwebsite.net/het/essays/uncert/subjective.html
}.
En este sentido, SAVAGE plantea que:
\begin{lnum}
    \item Existe un conjunto finito de \textit{estados de la naturaleza}\footnote{%
        Es interesante la definición que hace ARROW (1971) sobre el concepto de \textit{estados de la naturaleza}: Un estado es la resolución de la incertidumbre, “Una descripción del mundo tan completa que, si es cierta y conocida, las consecuencias que tiene aparejada cada acción serían también conocidas”. Implicito a esta definición es la noción de que solo \textit{un estado es verdaderamente cierto.}
    } cuyas probabilidades son asignadas de manera subjetiva por los decisores (permitiendo el modelo de inferencia);
    \item Los subconjuntos de estos \textit{estados} serán los \textit{eventos}, cuyas \textit{consecuencias} son descripciones de todo lo que pueda sucederle al decisor\footnote{%
        El conjunto de consecuencias es completamente arbitrario (en función del problema).
    };
    \item Una combinación de un acto del decisor y un \textit{estado de la naturaleza}, “seleccionado” por la naturaleza, determina una \textit{consecuencia}.
\end{lnum}
\textbf{[Definir el “conjunto de elección de altrnativas de riesgo”: conjunto de todas las colecciones o S-tuplas de funciones de distribución o loterías según el estado]}

Sobre este conjunto se establece una relación de preferencias sobre las que se imponen los siguientes axiomas:
\begin{lnum}
    \item \textit{Racionalidad.} Completitud y transitividad.
    \item \textit{Deseabilidad.}  Monotonía en sentido fuerte/ Insaciabilidad/ Continuidad. 
    \item \textit{Independencia extendida} Versión extendida (a niveles superiores, i.e. \text{estados}) del axioma de independencia sobre alternativas, que impone consistencia en la ordenación de las preferencias sobre los estados. Las preferencias entre actos depende solamente de las consecuencias en los eventos en que los dos actos que se estén comparando tengan valores diferentes. Axioma de \textbf{independencia de eventos disjuntos}. \textit{Estado-dependientes} (i.e. la función de utilidad sobre la que se valoran los distintos resultados puede cambiar con los resultados). 
    \item \textit{Independencia ordinal de los eventos.} La ordenación de preferencias entre consecuencias dentro de cada estado no depende de cuál sea el evento relevante que consideremos. Si dentro de cada estado el individuo prefiere una consecuencia a otra, eso es así independientemente del estado que da lugar a la comparación.
\end{lnum}
El modelo de la Utilidad Subjetiva de SAVAGE (1954) postula entonces una estructura de preferencias que permite: 
\begin{lnum}
    \item La expresión numerica de la valoración del decisor sobre las consecuencias, mediante una Función de Utilidad;
    \item La expresión numérica de las creencias del decisor sobre la probabilidad de los eventos, mediante una medida de probabilidad (aditiva de manera finita);
    \item La valoración de los actos mediante Esperanzas matemáticas de la Función de Utilidad de las consecuencias asociadas, con respecto a las probabilidades subjetivas de los eventos en los que dichas consecuencias se materializan;
    \item 
\end{lnum}
\textbf{[La genialidad del modelo se desprende de su simplicidad]}

\textbf{[Podemos plantear un modelo muy sencillo en los que el conjunto de actos del decisor sea: {llevar paraguas, no llevar nada} y el conjunto de estados de naturaleza sea: {llueve mucho, llueve poco, no llueve}. Su función de utilidad (discreta) puede expresarse de tal forma que sea “cóncava” respecto al eje de ordenadas (nivel de seco) y su utilidad creciente. Luego asigna una distribución de probabilidad cada vez que vaya a salir de casa: “Viendo el cielo”.]}

\textbf{[Plantear el modelo básico de SAVAGE]}
\subsubsection{Extensiones}
Posteriormente, ANSCOMBE y AUMANN (1963), ampliaron el modelo de SAVAGE (1954) con el propósito de incluir una “distribución de probabilidad objetiva”, exógenamente dada y común para todos los individuos (asociada a cada estado de naturaleza) sobre la cual el decisor lleva a cabo un acto/elección (de uno de los estados) en base a la distribución de probabilidad “subjetiva” que éste asigna a los eventos asociados a cada uno de estos estados\footnote{%
    Siguiendo la definición de ANSCOMBE y AUMANN (1963), las distribuciones de probabilidad de cada \textit{estado de la naturaleza}  representan \textit{loterías de ruleta} (con distribución de probabilidad conocida), mientras que las variables aleatorías (asociadas a cada \textit{estado de la naturaleza}) responden a \textit{loterías de caballo} (carecen de probabilidad objetivas). Por lo tanto, en este modelo una decisión es una apuesta en una carrera de caballos cuyos premios son apuestas en una ruleta de casino.
}.

\textbf{[Extender la versión de Savage al Modelo mentado, no hacer más]}
En consecuencia, la teoría de la probabilidad subjetiva puede verse como un caso más general de la teoría de la utilidad esperada, contemplando las situaciones de riesgo como un caso particular de las situaciones bajo incertidumbre.
\textbf{La gran ventaja del enfoque de la probabilidad subjetiva es la unificación de los problemas de riesgo y de incertidumbre bajo un mismo aparato formal, lo que implica un elevado grado de coherencia entre ambos.}

Sin embargo, plantea el inconveniente de que hereda las mismas del enfoque VON-NEUMANN-MORTGENSTERN que ahora veremos, así como otros fallos propios de esta especificación.
\subsection{Limitaciones y críticas a la TUE y a su extensión en forma de Probabilidad Subjetiva}
Los investigadores, a través de experimentos que permitiesen observar las elecciones reales de los individuos ante situaciones de riesgo, fueron descubriendo diferentes tipos de violaciones sistemáticas y recurrentes del modelo de utilidad esperada y sus hipótesis. Siguiendo a MACHINA (2008):
\subsubsection{Violaciones del axioma de independencia}
La violación más conocida es la paradoja de ALLAIS (1953), fruto de un experimento donde se pedía a los individuos que eligieran entre 2 pares de loterías, y en el que se observaron elecciones que contravenían sistemáticamente la Teoría de la Utilidad Esperada (https://youtu.be/FMEApkB4vUo).

En concreto, ALLAIS demostró que los individuos rara vez toman decisiones racionales cuando son sometidas a problemas de elección sucesivos inmediatamente. Los resultados del experimento de ALLAIS contravienen el axioma de independencia pues demuestra que los individuos cambian sus preferencias cuando las loterías se alteran en proporciones equivalentes. Particularmente, ALLAIS observó que ese sesgo hacía que los individuos tendieran a preferir loterías con pagos ciertos (\textit{efecto certeza}), es decir, loterías de varianza-cero. Podemos ilustrar la paradoja de Allais con el siguiente ejemplo:

\imagenfit{img1_P33.png}{[NO TITLE REQUIRED]}{[NO SOURCE REQUIRED]}

\footnote{%
    Al hacer frente al Juego 1, la mayoría de las personas elige la lotería A (ya que ofrece un mayor pago esperado y consideran las diferencias en probabilidad de obtener un pago positivo poco significativas). Sin embargo, en el \textit{Juego 2} (que es una trasnformación lineal de las loterías A y B) la mayoría de las personas prefiere el pago seguro.
}

Su investigación posterior permitió descubrir que esta paradoja no es sino un caso particular de dos efectos: \textit{el efecto consecuencia común}\footnote{%
    Se denomina así porque las cuatro loterías que se necesitan se forman sobre un conjunto de resultados o consecuencias comunes. Por ejemplo, genera mucha más utilidad recibir la cantidad “X” como premio máximo de una lotería que recibirla como premio mínimo de otra lotería.
} {y el efecto razón común}\footnote{%
    Se trata de un experimento parecido al anterior, con elección entre dos pares de loterías, pero en este caso la razón o cociente entre probabilidades de los resultados o pagos no nulos se mantiene constante. Apunta a que la gente sufre cierta modificación de sus preferencias por el riesgo cuando son confrontados a elecciones sobre loterías que guardan cierta relación de equivalencia entre ellas. Una crítica de esta evidencia (MACHINA, 2008) es que en diversos experimentos se ha observado que los sujetos típicamente corrigen su comportamiento una vez se les revela la naturaleza de sus violaciones de la utilidad esperada.
}.
\subsubsection{Violaciones del supuesto de invarianza descriptiva y procedimental}
(En este tipo de estudios cabe destacar la contribución de RICHARD THALER\footnote{%
    RICHARD THALER fue galardonado con el Premio Nobel de Economía en 2017 «\textit{Por sus contribuciones a la economía conductual».}
}):

Por un lado, se han encontrado violaciones sistemáticas del supuesto subyacente estándar de estabilidad de las preferencias y de invarianza con respecto a la descripción del problema en elecciones de riesgo, de modo que variaciones en el “\textit{marco}” o en la forma (\textit{frame}) de presentar un mismo problema conllevan elecciones diferentes por parte de los sujetos. Algunos ejemplos de este tipo de \textit{framing effects} son:
\begin{lnum}
    \item 1.	El fenómeno del \textit{“punto de referencia”}, pues las actitudes de riesgo hacia ganancias y pérdidas no pueden ser explicadas satisfactoriamente con una rígida función de utilidad sobre los niveles de riqueza total final sino que podrían explicarse mejor mediante funciones definidas en términos de cambios o variaciones desde el “punto de referencia” de la riqueza presente, siguiendo la idea de MARKOWITZ.
    \item La \textit{descripción del problema}, según este se presente en forma de matriz, árbol de decisión, ruedas de ruleta, enunciados escritos.
    \item El \textit{enunciado} de un mismo problema de riesgo en forma de decisión de apostar frente a la decisión de asegurarse.
\end{lnum}
Por otro lado, se ha encontrado evidencia de que formatos de respuesta alternativos\footnote{%
    Por ejemplo, pidiendo que los sujetos expresen sus preferencias sobre unas loterías especificando sus equivalentes de certeza, de ganancia o de probabilidad.
} conducen a ordenaciones diferentes, y, por ende, inconsistentes (\textit{response-mode effects}).
\subsubsection{Paradoja de ELLSBERG}
Mencionadas las críticas generales asociadas al marco conceptual de la Teoría de la Utilidad Esperada, cabe destacar una crítica adicional dirigida específicamente contra el modelo de probabilidades subjetivas (o modelo genérico de incertidumbre) por ser especialmente relevantes las conclusiones extraídas de ésta.

Hablamos en concreto de la paradoja de ELLSBERG (1961), siendo ésta el ejemplo más conocido sobre las desviaciones empíricas e inconsistencias que presenta su aplicación real. ELLSBERG (1961) demostró\footnote{%
    Una variante simple de éste experimento consiste en que se da a elegir entre dos pares de apuestas: A) Una urna donde hay 100 bolas de dos colores diferentes y cuya proporción es conocida; B) Una urna donde hay 100 bolas de dos colores diferentes, pero cuya proporción es desconocida. El decisor debe escoger de cuál de las dos urnas desea extrae dos bolas, con el fin de que las extraídas sean del mismo color (ganar la apuesta).
} que las decisiones de elección violaban mayoritariamente la teoría de la probabilidad subjetiva, pues las probabilidades subjetivas que de las elecciones en incertidumbre se inferían no eran coherentes con la utilidad esperada. Lo que sí se observaba era una tendencia a preferir apuestas bajo riesgo frente a apuestas bajo incertidumbre, incluso aunque estas últimas fuesen aparentemente más ventajosas según la teoría de la probabilidad subjetiva. \textbf{Esto sugiere que los individuos interpretan la incertidumbre como un mal en sí mismo, como si generase rechazo per se. Es lo que se ha denominado “aversión a la ambigüedad”}\footnote{%
    A diferencia de los conceptos económicos de \textit{riesgo y aversión al riesgo}, no existe un acuerdo unánime sobre lo que \textit{aversión a la ambigüedad} e incluso \textit{ambigüedad} misma son exactamente.
}.

\imagenfit{img1_P35.png}{[NO TITLE REQUIRED]}{[NO SOURCE REQUIRED]}

Algunas de las deficiencias que presentaba la teoría de la probabilidad subjetiva (comunes también en el modelo de riesgo) pueden ser solventados mediante un modelo de utilidad no esperada que sea \textit{probabilísticamente sofisticado}\footnote{%
    Siguiendo la axiomatización de MACHINA y SCHMEIDLER (1992), la \textit{sofisticación probabilística} significa que el agente puede estabecer algún tipo de probabilidad subjetiva sobre el espacio de sucesos, lo que permite luego proceder como en un modelo de utilidad no esperada para probabilidad objetivas. El modelo de utilidad esperada más popular en este tipo de situaciones es el \textit{rank-dependant utility} en incertidumbre (SCHMEIDLER, 1989). La evidencia empírica parece apuntar que las funciones de ponderación $W$ utilizadas (Ver Anexo) para incertidumbre exhiben patrones similares a los encontrados en riesgo. Igualmente, y de manera análoga al caso de riesgo, se encuentra \textit{“orientación hacia los extremos”} (sobre ponderación de los mejores y peores resultados) y ausencia de sensibilidad hacia los resultados intermedios. La sofisticación probabilística permite acomodar los fenómenos derivados de la paradoja de Allais (desviación de linealidad de las probabilidades), pero no las desviaciones de la probabilidad subjetiva puestas de manifiesto por la paradoja de Ellsberg. En efecto, la paradoja de Ellsberg viola la sofisticación probabilística:\textit{ se demuestra que la aversión a la ambigüedad implica que la función de ponderación (W) debe ser convexa (en su definición específica para conjuntos), lo cual no es posible con sofisticación probabilística}. Existen otros modelos de incertidumbre distintos de los orden-dependientes, como los modelos de múltiples-previos (\textit{multiple priors}, en especial: WALD, 1950; ARROW y HARWICZ, 1972; SAVAGE, 1954; o, muy utilizado actualmente, GILBOA y SCHMEIDLER, 1972), los enfoques de ambigüedad libres de modelo (\textit{model-free}: EPSTEIN, 1999) y los de múltiples etapas (\textit{multi-stage}: ERGIL y GUL, 2004).
}.

En cualquier caso, si bien la Teoría de la Utilidad Esperada (y Subjetiva) no es capaz en muchas ocasiones de predecir qué decisiones tomarán agentes reales concretos, sí constituye una referencia básica a la hora de determinar qué decisión deberán tomar. De esta forma, adquiere una \textit{dimensión normativa} que sobrepasa la dimensión positiva inicial\footnote{%
    Parece razonable pensar que en la medida en que los agentes estén plenamente informados y estén dispuestos a asumir el coste de utilizar esa información para tomar decisiones racionales, se comportarán siguiendo en muchos casos los dictados de la Teoría de la Utilidad Esperada (y/o, Subjetiva).
}. En definitiva, las regularidades empíricas que acabamos de mencionar y que contradicen ciertos supuestos del enfoque de la Teoría de la Utilidad Esperada han guiado el desarrollo de enfoques teóricos alternativos posteriores que permitan (en cierto grado), acomodarlas.
\clearpage
\section{Críticas y alternativas a la teoría de la utilidad esperada}
Conocidos como Teorías de la Utilidad No Esperada (non-expected utility theories)\footnote{Un breve resumen de los diferentes modelos se puede ver en el [Anexo].}, entre los quedebemos destacar la \textit{Prospect Theory}.
\subsection{Origenes de la \textit{Prospect Theory}}
Fruto de la dilatada colaboración entre dos psicólogos, DANIEL KAHNEMAN y AMOS TVERSKY\footnote{%
    Que se inició en los años 50 en el ejército israelí, donde el primero trabajaba como psicólogo y el segundo como capitán de paracaidistas.
}, el genesis de esta Teoría se remonta a marzo de 1979, cuando ambos la enunciaron (en su primera versión) en un artículo publicado en la revista Econometrica\footnote{Puede verse en: https://www.jstor.org/stable/1914185.}.

Esta primera versión presentaba algunas limitaciones\footnote{%
    La limitación más relevante era el no cumplimiento de la dominancia estocástica de primer orden, es decir, que los agentes elegirían en ocasiones loterías dominadas.
}, los autores mejoraron su análisis mediante la incorporación de probabilidades acumuladas y publicaron una revisión en el año 1992\footnote{Puede verse en:\\ https://link.springer.com/article/10.1007/BF00122574.}.

Según señala PETER BERNSTEIN en su entretenida explicación de la teoría, KAHNEMAN y TVERSKY eligieron ese nombre porque era llamativo y memorable, no porque guardara relación con su contenido.

La Prospect Theory tiene base empírica y aspira a reflejar cómo la gente se comporta en realidad, no cómo debiera hacerlo si fuera racional. No es, pues, una teoría normativa, sino empírica y positiva.
\subsection{Diferencias con la TUE}
Se distancia de la TUE en lo que se refiere a tres grandes cuestiones:
\begin{lnum}
    \item La definición de las alternativas sobre las que versan nuestras decisiones humanas;
    \item La valoración que les damos; y
    \item La ponderación que, a la vista de su probabilidad, les atribuimos.
\end{lnum}
\subsection{Definición de las alternativas: reglas heurísticas}\footnote{%
    Las reglas heurísticas son atajos mentales o reglas simplificadas que las personas usan para tomar decisiones complejas de forma rápida e intuitiva, sin hacer un análisis exhaustivo de todas las probabilidad y resultados (este punto es el que más marca distancia de la Teoría de la Utilidad Esperada). KAHNEMAN y TVERSKY (1974) enfatizan que el comportamiento humano no es irracional, sino adaptativo: las personas usan reglas simples que funcionan bien la mayoría del tiempo, pero producen sesgos sistemáticos en contextos complejos. Por tanto, en la Prospect Theory, las reglas heurísticas sustuyen a los procesos racionales de la Utilidad Esperada o de la Utilidad Subjetiva, siendo mecanismos congnitivos que permiten decidir con esfuerzo mínimo aunque a veces conduzcan a sesgos sistémicos. En “Judgement under Uncertainty: Heuristics and Biases” (KAHNEMAN y TVERSKY, 1974: https://sites.socsci.uci.edu/~bskyrms/bio/readings/tversky_k_heuristics_biases.pdf) identifican tres grandes heurísticas que la gente usa para juzgar probabilidades y riesgos: \\ TABLA P37
}

La Prospect Theory parte de que nuestra limitada capacidad intelectiva nos obliga a simplificar los problemas de decisión que se nos plantean, y lo hacemos siguiendo ciertas reglas (heuristic rules):
\begin{itemize}
    \item •	La primera y principal es que (recuperando a MARKOWITZ), al enjuiciar alternativas, comparamos no valores absolutos –como supone la TUE–, sino variaciones o cambios respecto a cierto nivel que tomamos como punto de referencia. Así pues, las alternativas las vemos en términos de ganancias o pérdidas respecto a cierto nivel de referencia. Ese nivel de referencia suele ser el statu quo, pero puede ser también el nivel psicológico al que aspiramos o incluso algún nivel arbitrario que, sin darnos cuenta, nos ha sugerido aquél que nos ha planteado la cuestión.
    \item La segunda regla, relevante cuando nos enfrentamos a una serie de acontecimientos, se refiere a su combinación. Así, por ejemplo, una ganancia seguida de una pérdida más pequeña, \textit{¿las percibiremos psicológicamente como una ganancia neta o, por el contrario, las mantendremos intelectualmente separadas y las percibiremos de forma separada?}\footnote{%
        Al igual que ocurre con las reglas sobre acumulación de rentas o declaración separada en un IRPF progresivo, esas reglas de combinación o acumulación de sucesos tienen gran importancia.
    }
    \item Finalmente, una de las predicciones de la Prospect Theory es que, aunque no se modifique el fondo de las alternativas, un cambio en el marco de referencia puede alterar nuestra elección (\textit{framework effect}), pues nos sentimos atraídos por las ganancias ciertas y rehuimos las pérdidas seguras\footnote{%
        Ese \textit{framework-effect} queda ilustrado en un clásico experimento en dos etapas:
    \begin{lnum}
        \item Le regalamos a un sujeto 1.000 € y le decimos que en una segunda fase tiene que elegir entre dos premios adicionales: a) 500 € más, seguros; o b) 1.000 € más, pero con una probabilidad del 50 \%.\\
        Casi todo el mundo prefiere los 500 € adicionales, asegurándose 1.500 € en total. 
        \item En un segundo experimento, regalamos a un sujeto 2.000 € y le decimos que, en una segunda fase, tiene que escoger entre dos multas, que restarán de sus 2.000 €: a) perder 500 € seguro; o b) perder 1.000 € con una probabilidad del 50\%.\\
        Casi todo el mundo se inclina ahora por la segunda alternativa, porque se resiste a perder con seguridad 500 € de sus 2.000 €.
    \end{lnum}
    Ahora bien, ese par de elecciones resulta paradójico, porque en ambos casos se está dando a elegir entre lo mismo: a) 1.500 € seguros; o b) 1.000 € ó 2.000 € con probabilidad del 50\%. Como en el primero se utiliza un marco de referencia de ganancias, la gente suele ser menos arriesgada. Sin embargo, en el segundo experimento el marco de referencia hace alusión a pérdidas, por lo que la gente suele estar más dispuesta a jugársela.
    }.
\end{itemize}
\subsubsection{Valoración de las alternativas: la función de valor en forma de S}
Al igual que la TUE atribuye a cada resultado cierta utilidad, la \textit{Prospect Theory} atribuye a cada alternativa (entendida como ganancia o pérdida respecto a nuestro nivel de referencia) un cierto valor. Esa función de valor tiene la forma aproximada de S que aparece en el siguiente gráfico:

\imagenfit{img1_p38.png}{[NO TITLE REQUIRED]}{[NO SOURCE REQUIRED]}

Esta función de valor se caracteriza por lo siguiente:
\begin{itemize}
    \item •	En materia de ganancias: el valor marginal –el que atribuimos a cada nueva unidad– es cada vez menor (esto es, la curva va perdiendo pendiente). Por eso, en materia de ganancias somos conformistas, y preferimos una ganancia cierta a otra mayor pero hipotética\footnote{%
        Como bien se expresa en el Quijote: \textit{“Más vale pájaro en mano que buitre volando”}; otra referencia: \textit{“No dejes lo ganado por lo que has de ganar”} (Libro de Buen Amor).
    }. Esto implicará que la función de valor es cóncava en el lado de las ganancias.
    \item En materia de pérdidas ocurre algo parecido: su impacto marginal es cada vez menor. Por eso no nos importa arriesgarnos a sufrir grandes pérdidas si con ello evitamos una pérdida menor pero cierta\footnote{%
        \textit{“De perdidos, al río”}. Ante una situación muy difícil, se opta por la solución más descabellada. Tal opción se debate ante la desesperación que conduce a que ya nada importe, o ante la idea de que lo más absurdo representa la única solución.
    }. Así pues, a diferencia de la TUE, la \textit{Prospect Theory} pronostica que, en materia de pérdidas, somos amantes del riesgo (\textit{risk-seekers}). Esto implicará que la función de valor es convexa en el lado de las pérdidas.
    \item •	En las inmediaciones del origen de coordenadas, la pendiente de la curva en el tramo de pérdidas (esto es, en el cuadrante inferior izquierdo) es mucho mayor que en el de ganancias (cuadrante superior derecho) \footnote{%
        \textit{“De perdidos, al río”}. Ante una situación muy difícil, se opta por la solución más descabellada. Tal opción se debate ante la desesperación que conduce a que ya nada importe, o ante la idea de que lo más absurdo representa la única solución.
    }. Así pues, la S es asimétrica, y su tramo descendente es más vertical que el ascendente. Esa asimetría refleja la aversión a las pérdidas\footnote{%
        Siempre rechazamos una apuesta que nos ofrezca ganar o perder la misma cantidad con una probabilidad del 50 \%, pues las pérdidas nos duelen más de lo que nos alegran las ganancias de igual importe. Una manifestación directa de lo anterior es el llamado \textit{endowment effect} (efecto dotación): en general, pedimos mucho más por desprendernos de algo que ya tenemos (pérdida) que lo que estaríamos dispuestos a pagar por adquirirlo (ganancia).
    }.
\end{itemize}
\subsubsection{Ponderación de las alternativas: pesos decisorios y efecto certeza}
Al igual que la TUE, los valores o utilidades atribuidos a cada alternativa se ponderan por cierto peso antes de su comparación definitiva. Sin embargo, con el fin de reflejar la paradoja de Allais y nuestra hipersensibilidad a pequeños riesgos, la \textit{Prospect Theory} no pondera las alternativas con sus probabilidades objetivas, sino con ciertos pesos decisorios \textit{(decision weights)} que guardan una relación no lineal con las probabilidades\footnote{%
    Son mayores que estas cuando las probabilidades son bajas. Luego pierden sensibilidad (pendiente) en los tramos centrales de probabilidad y la recuperan de nuevo para probabilidades muy altas, atraídas por el imán del efecto certeza (tal y como sugiere la paradoja de Allais).
}. Esto se muestra en el siguiente gráfico: 

\imagenfit{img1_p40.png}{[NO TITLE REQUIRED]}{[NO SOURCE REQUIRED]}

Ponderacion de las alternativas

\subsection{Coherencia con la evidencia empírica: \textit{endowment effect} y paradoja de EASTERLIN}
Esta teoría da una explicación coherente a muchos fenómenos económicos y políticos, difícilmente reconciliables con la TUE y la racionalidad de los agentes económicos. Algunos de los más destacados son:
\begin{lnum}
    \item 1.	Las ventajas de las devaluaciones y la inflación para alterar los precios relativos y salarios reales: los salarios y precios nominales son inflexibles a la baja, porque su reducción se ve como una pérdida; por el contrario, su falta de ajuste a la inflación se tiende a ver como la renuncia a una ganancia, lo cual es menos doloroso.
    \item En las épocas de recesión, el mercado de viviendas de segunda mano tiende a paralizarse, ante la inflexibilidad a la baja de los precios de venta. Esa inflexibilidad puede atribuirse a la renuncia de los propietarios a incurrir en pérdidas respecto a su nivel de referencia (\textit{endowment effect}).
    \item La reacción social ante riesgos que se consideran nuevos es muy virulenta, aunque los riesgos sean muy remotos (\textit{Vacas Locas}, 1996; Organismos Genéticamente Modificados, 1997, etc.). De forma parecida por eliminar todo vestigio de ciertos riesgos (en Estados Unidos, limpieza de espacios medioambientales) se gastan fortunas que no se aplican a reducir riesgos inevitables, pero mucho mayores (por ejemplo, accidentes en carretera). Así pues, el peso decisorio que nuestra mente atribuye a riesgos que consideramos nuevos es muy superior al de su probabilidad objetiva.
    \item La paradoja de Easterlin: Según constató en los años setenta RICHARD EASTERLIN (y se ha corroborado posteriormente varias veces), aunque aumente la renta por habitante de un país, el porcentaje de sus habitantes que se declaran felices suele permanecer estable. La razón es que la felicidad es un concepto relativo, para el que tomamos como referencia el nivel de bienestar de nuestros vecinos\footnote{%
        Parece una variante de la Teoría de la identidad social (TAJFEL y TURNER, 1979), cuya idea central es que las personas no tienen una identidad única y fija, sino múltiples identidades sociales y la identidad que se activa depende del contexto en el que el individuo se encuentra \textit{(identity salience)}. Otra proupesta formal es la \text{Self-categorization Theory }(TURNER, 1987), que complementa a la anterior en tanto que: 1) el individuo categoriza el entorno en términos de grupos \textit{“nosotros”} (ingroup) y \textit{“ellos”} (outgroup); 2) la categoría más diagnóstica o distintiva en una determinada situación se vuelve saliente. En economúa del comportamiento, autores como AKERLOF y KRANTON (2000) incorporan este fenómeno al análisis económico. Éstos llaman “identidad”a esa autoimagen contextual, y formalizan cómo las personas eligen comportamientos que maximicen su utilidad de identidad, no solo su utilidad material. Ver: https://www.researchgate.net/publication/24091706\_Economics\_And\_Identity
    }.
\end{lnum}
\subsection{Valoración}
A pesar de que tienen tan solo cuatro décadas de vida la \textit{Prospect Theory} (1979) y la \textit{Psicología de las Finanzas}\note{Aplicación al \texit{Behavioral Finance} [ver tema 3.B.24]\\
Además dee resultar coherente con los fenómenos descritos, la Prospect Theory y la Psicología de las Financas explican varios enigmas financieros o anomalías difíciles de explicar desde una óptica tradicional:
\begin{lnum}
    \item \textit{Equity premium puzzle}: ¿Cómo explicar que, según reiterados cálculos, la inversión en acciones produzca, de forma sistemática a largo plazo una rentabilidad superior en más de 6 puntos porcentuales a la de la deuda pública? La \textit{Prospect Theory} lo achaca a la gran frecuencia con la que los inversores evalúan el valor de sus carteras y a la asimetría con la que reaccionamos a las pérdidas y ganancias. Esa evaluación frecuente hace muy dolorosas las pérdidas que provocan las fluctuaciones a corto plazo de las cotizaciones incluso aunque luego sean compensadas por ganancias. Tal aversión miope a las pérdidas hace que, en equilibrio, los inversores exijan de las acciones un rendimiento superior al de la renta fija. Un corolario –refrendado experimentalmente– es que la aversión por la renta variable disminuye cuanto más se alarga el período de valoración de las carteras. 
    \item \textit{Volatility puzzle}: ¿Por qué las cotizaciones de las acciones –que, en teoría, reflejan tan sólo el valor presente de su flujo futuro de dividendos– oscilan mucho más que los propios dividendos? La Prospect Theory lo fundamenta en el proceso psicológico de combinación de ganancias y pérdidas. En época de bonanza, cuando los dividendos y la Bolsa suben, los inversores acumulan ganancias latentes y pierden miedo a las futuras pérdidas, pues si éstas se producen tan sólo minorarán las ganancias: no se verán, pues, como genuinas pérdidas, sino como menores ganancias. Ese neteo psicológico reducirá su aversión al riesgo y su tasa subjetiva de descuento, lo que los llevará a estar dispuestos a pagar más por acción para un mismo nivel de dividendos. Así pues, un aumento de dividendos, al disminuir la prima de riesgo exigida por los inversores, hará que las cotizaciones suban por doble motivo. El proceso contrario se dará cuando bajen los dividendos. En suma, el efecto colchón de las ganancias acumuladas durante la bonanza (llamado house money effect, en alusión al optimismo típico del jugador de casino que va ganando) tendrá un efecto procíclico.
    \item \textit{Disposition effect}: ¿Por qué los inversores son reacios a vender acciones con pérdida y proclives a vender las que arrojan plusvalías? La Prospect Theory lo atribuye a la forma de S de la función de valor, que refleja el valor marginal decreciente de ganancias y pérdidas. Imaginemos, por ejemplo, que un inversor compró una acción en 50 €. La acción ahora cotiza a 55 € y existe la misma probabilidad de que suba a 60 € o de que baje a 50 €. Pues bien, el valor atribuido a una plusvalía segura de cinco euros será mayor que el de una potencial ganancia de 10 € con una probabilidad del 50\%. El inversor materializará pues su ganancia. Imaginemos ahora que esa misma acción cotice a 45 € –con una minusvalía latente de cinco– y que existe la misma probabilidad de que suba a 50 € o de que baje a 40 €. Si vende, el inversor materializará una pérdida de 5 €. Si no lo hace, su potencial minusvalía será de 10 € con una probabilidad del 50\%. Enfrentando a esa desagradable tesitura, el dolor marginal decreciente de las pérdidas hará que prefiera conservar el valor.
    \textit{Dividend puzzle}: ¿Por qué en Estados Unidos muchas empresas han seguido pagando dividendos, a pesar de que –por lo menos hasta ahora– tributaban más que las plusvalías? La Prospect Theory supone que las empresas razonan como Maquiavelo: \textit{por la misma razón que los actos de severidad deben hacerse de una vez para que dejando menos tiempo para notarlos ofenderán menos, los beneficios deben otorgarse poco a poco, a fin de que puedan saborearse mejor}. De forma parecida, si una empresa consigue unos beneficios de 10 € por acción, podrá beneficiar al accionista si, en vez de transmitírselos en bloque, los fracciona (por ejemplo, en un dividendo de 2 € y una plusvalía latente de 8 €). Esa segregación hará que el accionista se beneficie de dos ganancias separadas de dos y de ocho, que disfrutará más que una sola de 10 €. Por un motivo parecido, a una empresa con pérdidas puede interesarle repartir un pequeño dividendo, porque dado el valor marginal decreciente que atribuye a aquéllas, el accionista preferirá un dividendo de dos y unas pérdidas brutas de 12 € a unas pérdidas netas de 10 €. En suma, las reglas psicológicas de combinación y segregación de sucesos harán que reaccionemos de forma distinta ante lo que, desde el punto de vista financiero, es idéntico.
\end{lnum}
} se han aplicado ya a muchos otros campos\footnote{%
    Por ejemplo, destaca su creciente aplicación en el mundo jurídico sobre todo para explicar sesgos típicos en las decisiones de los jurados.
}.
Representan un enfoque prometedor, complementario del de aquellas concepciones que suponen la perfecta racionalidad de las personas hacen comprensible lo que hasta ahora tan solo parecía \textit{anómalo}\footnote{%
    Representan, en fin, para el mundo de la economía y las finanzas, algo parecido a lo que “El príncipe” de MAQUIAVELO o, más recientemente, “Una teoría económica de la democracia” de ANTHONY DOWNS representaron para el de la política..
}.
{\color{red} 
https://www.nobelprize.org/prizes/economic-sciences/2002/kahneman/lecture/ \\
 Bounded rationality – Tiene buena pinta \\
https://www.nobelprize.org/prizes/economic-sciences/2017/thaler/lecture/ \\
Leer los pdfs de ambos.}
\clearpage
\section{Conclusion}
\subsection{Recapitulación (Ideas clave)}
En la presente exposición hemos analizado la TUE. Se trata de una contribución realizada por una serie de economistas que aplican la teoría de la elección a situaciones de ausencia de información perfecta.

Hemos abordado la axiomática, los componentes de la función de utilidad esperada, las aplicaciones, las críticas y el enfoque de probabilidades subjetivas.

Estas contribuciones son de utilidad en varias ramas de la teoría económica.

\subsection{Relevancia}
[¿Completar?]

\subsection{Extensiones y relación con otras partes del temario}
Los desarrollos aquí expuestos tienen aplicaciones en otras ramas de la teoría económica.\\ Por ejemplo:\\
Teoría del equilibrio general en condiciones de incertidumbre: ARROW y DEBREU (1953, 1959) incorporan el enfoque de preferencia por el estado y un contexto dinámico49.

[La clave de los mercados Arrow-Debreu es que en t=0 hay mercados para todos los bienes en cada uno de los estados. Una idea propuesta por ARROW es que es posible reducir los mercados en t=0 a un único bien contingente, a cambio de que en t=1 se reabran los mercados de contado en todas las mercancías. El equilibrio de Radner es para el equilibrio en el mercado de activos.]

Economía financiera: En especial en la teoría de la selección de carteras [ver tema 3.B.23]. El modelo de MARKOWITZ parte de la teoría de la utilidad esperada para estudiar la elección del consumidor entre activos financieros. El mismo sustento teórico es empleado por TOBIN para extraer la función de demanda de dinero (entendiendo el dinero como un activo).

Además, las críticas a la TUE también se han incorporado a las economías financieras dentro del campo del behavioral finance.

Teoría de la demanda de consumo de HALL plantea un modelo en un entorno estocástico y muestra que con una función de utilidad cuadrática, el consumo en un período sigue un paseo aleatorio [ver tema 3.A.33].
\subsection{Opinión}
[Completar]
\subsection{Idea final (Salida o cierre)}
En definitiva, conocer las raíces teóricas de la elección del consumidor bajo riesgo e incertidumbre es fundamental a la hora de comprender desarrollos en otros campos de la teoría económica que son conocedores de que la incertidumbre es un aspecto ineludible a la vida humana.

\clearpage
\pagenumbering{gobble}
\MyBibliography
\MyBibItem{Autor}{Título}{Año}{DOI -- LINK}


%anexos
\clearpage
\anexo{\textbf{Preguntas Test}}
%\input{\TemaCode/questions.tex} 



\end{document}